\chapter{Analisi dello stato del arte e contesto tecnologico} %\label{1cap:spinta_laterale}
% [Analisi dello stato del arte] {Analisi dello stato del arte e contesto tecnologico}
%

%\begin{citazione}
%In questo capitolo analizzeremo il contesto tecnologico su cui siamo andati a lavora.
%\end{citazione}

\section{OWASP Cloud Security Top 10}

Il portale AppSecMaster ha elaborato una propria classificazione dei dieci rischi più comuni e di maggiore impatto nei moderni ambienti cloud, ispirandosi ai principi promossi da OWASP in materia di sicurezza applicativa, che vanno dall'infrastruttura mal configurata alla scarsa visibilità delle risorse in esecuzione. Di seguito è riportata la lista delle principali problematiche individuate:
\begin{enumerate}
    \item \textbf{Security Misconfiguration} — Bucket di storage lasciati con permessi troppo permissivi o istanze di calcolo distribuite senza le dovute restrizioni restano tra le cause più frequenti di violazione dei dati in ambienti cloud.
    \item \textbf{Identity / Access Mismanagement} — Ruoli IAM assegnati con privilegi superiori a quelli effettivamente necessari, e mai rivisti dopo la configurazione iniziale, espongono i sistemi a un concreto rischio di escalation dei permessi.
    \item \textbf{Improper Secrets/Key Management} - Se chiavi API e credenziali vengono salvate in chiaro o incorporate nel codice sorgente, la loro compromissione può aprire un accesso laterale a più servizi collegati.
    \item \textbf{Insecure Network Segmentation} — Senza una reale separazione tra reti e ambienti diversi, un attaccante che viola un singolo componente può muoversi liberamente verso risorse più sensibili.
    \item \textbf{Injection / Input Validation Flaws in Cloud Services} — Le classiche vulnerabilità di injection assumono forme nuove nei contesti cloud-native, colpendo funzioni serverless, code di messaggi e gateway API dove i controlli sugli input non sono sempre applicati in modo uniforme.
    \item \textbf{Insecure Deserialization / API Exposure} — La gestione non sicura di dati serializzati, unita ad API interne prive di adeguati controlli di autenticazione o versionamento, costituisce un punto d'ingresso sfruttabile per eseguire codice non autorizzato.
    \item \textbf{Vulnerable Third-Party/Container Dependencies} — Librerie open source e immagini container non aggiornate accumulano vulnerabilità note, soprattutto quando le stesse basi vengono riutilizzate tra più progetti senza verifiche periodiche.
    \item \textbf{Improper Asset Inventory / Visibility} — Risorse cloud dimenticate o mai tracciate correttamente, incluse istanze abbandonate e servizi "ombra", sfuggono facilmente al controllo dei team di sicurezza.
    \item \textbf{Insufficient Logging / Monitoring} — Log distribuiti su decine di servizi diversi e privi di un formato coerente rendono difficile individuare per tempo un'intrusione in corso.
    \item \textbf{Unvalidated Redirects/Cloud Function Abuse} — Funzioni serverless eseguite con permessi troppo ampi, combinate a controlli insufficienti sui redirect, possono essere sfruttate per aggirare le restrizioni di accesso e ottenere un punto d'appoggio nascosto nell'infrastruttura.


    
\end{enumerate}

Tra le problematiche precedentemente descritte , il presente lavoro si basa in particolarme su Insecure Cloud Configurations e Insecure API Exposure per via della loro rilevanza diretta rispetto agli obiettivi del sistema sviluppato e alla loro frequente ricorrenza negli scenari di autorizzazione tipici delle architetture basate su API.

\subsection{Security Misconfiguration}
Una misconfiguration, in ambito cloud, nasce da un errore nella configurazione
di un servizio, di un prodotto o di una risorsa infrastrutturale: le
impostazioni di sicurezza predefinite non vengono irrobustite a
sufficienza, e questo lascia scoperto un fianco che con poco sforzo potrebbe
essere protetto. I numeri parlano chiaro: il documento ufficiale OWASP
riporta che il 90\% delle applicazioni testate è risultato vulnerabile ad
almeno una forma di misconfiguration~\cite{owasp_a05_2021}. Nella maggior
parte dei casi non c'è nulla di sofisticato dietro: è una dimenticanza,
un'impostazione lasciata di default, un permesso che nessuno ha mai rivisto.
Ed è proprio questo che la distingue da una vulnerabilità nel codice
applicativo — qui il problema non sta nella logica del software, ma in come
le risorse sono state predisposte.

Le forme più comuni in cui una misconfiguration si manifesta sono, in fondo,
poche e ricorrenti. Un bucket di storage lasciato accessibile pubblicamente,
quando avrebbe dovuto restare privato, è probabilmente la più diffusa: basta
una policy impostata male, magari copiata da un altro progetto senza troppa
attenzione. Un container che gira con privilegi di root, invece, espone il
sistema a un rischio diverso ma altrettanto concreto — se l'attaccante riesce
a comprometterlo, si ritrova con permessi ben oltre quelli che gli
servirebbero per fare danni. C'è poi il caso dei template
Infrastructure-as-Code scritti senza le dovute verifiche: un file Terraform
o CloudFormation che definisce una risorsa con permessi troppo ampi si porta
dietro lo stesso errore ogni volta che viene riutilizzato, replicandolo su
ambienti diversi senza che nessuno se ne accorga.

\paragraph{Un caso reale: il data breach Capital One}
Il caso Capital One del luglio 2019 resta probabilmente l'esempio più citato
di questa categoria di rischio. Un bucket Amazon S3 lasciato con permessi
troppo ampi permise a un'attaccante di sottrarre oltre 100 milioni di record
clienti, tra dati anagrafici, indirizzi e informazioni finanziarie
sensibili~\cite{sentinelone_cloud_misconfig}. Nessuna falla nella piattaforma
cloud era coinvolta: la responsabilità ricadeva interamente sulla
configurazione impostata dal cliente. Non a caso AWS, commentando
l'accaduto, ribadì che \textit{"Amazon S3 is secure by
default"}~\cite{cloudcomputingnews_2023}, scaricando di fatto la colpa sulla
gestione dei permessi lato utente. Sul piano economico il conto fu salato:
80 milioni di dollari di multa dall'Office of the Comptroller of the
Currency, più 190 milioni versati per chiudere una class action. Capital One
dimostra bene quanto una semplice disattenzione in fase di configurazione
possa costare quanto, se non più di, un attacco costruito su un exploit
tecnicamente elaborato.

\bigskip

Episodi come questo spiegano perché, negli ultimi anni, le organizzazioni
abbiano iniziato ad affidarsi sempre meno alla revisione manuale delle
configurazioni e sempre più a strumenti automatizzati capaci di individuare
questo tipo di errori prima che vengano sfruttati. Una configurazione
insicura, del resto, raramente viene introdotta di proposito: emerge quasi
sempre da un template riutilizzato senza controllo, da un default mai
cambiato, da una revisione saltata per mancanza di tempo — fattori che un
controllo manuale sistematico difficilmente riesce a intercettare su larga
scala e con continuità.

Il documento ufficiale OWASP illustra questa categoria di rischio attraverso
quattro scenari esemplificativi~\cite{owasp_a05_2021}. Il primo riguarda
applicazioni di esempio dimenticate su un server di produzione, spesso
accompagnate da account amministrativi con credenziali mai modificate
rispetto ai valori predefiniti. Il secondo descrive un server su cui la
funzionalità di directory listing resta abilitata, consentendo a un
attaccante di risalire fino al codice sorgente tramite reverse engineering.
Il terzo riguarda messaggi di errore troppo dettagliati restituiti
dall'application server, capaci di esporre informazioni sull'infrastruttura
sottostante o versioni di componenti vulnerabili. Il quarto, infine,
riguarda permessi di condivisione predefiniti di un cloud service provider
lasciati aperti ad altri utenti su Internet, con conseguente esposizione di
dati sensibili archiviati nello storage cloud.

Fra i quattro scenari, due ricadono nel perimetro del sistema presentato in questo lavoro e due ne restano fuori.




\section{API Security Risks}

Le API oggi fanno gran parte del lavoro che una volta faceva il sito web
stesso: non restituiscono una pagina pensata per essere letta da una
persona, ma dati grezzi destinati a un'altra app, a un altro servizio, a un
sistema che poi decide cosa farne. Lo schema delle richieste è prevedibile,
gli endpoint espongono operazioni precise su risorse precise. Per un
attaccante questo significa meno lavoro di interpretazione: non deve capire
un'interfaccia grafica, gli basta osservare come sono costruite le chiamate.

OWASP mantiene infatti, accanto al Top 10 per le applicazioni web, una
classificazione separata dedicata alle API. Alcune categorie compaiono in
entrambe le liste — la misconfiguration, il broken access control — ma buona
parte dei rischi specifici delle API riguarda un livello che il Top 10
generale non copre: l'autorizzazione su singoli oggetti, proprietà,
funzioni. Un'applicazione web tradizionale filtra cosa mostrare all'utente;
una API deve filtrare cosa restituire nella risposta. La differenza sembra
minima, ma è da lì che nasce buona parte delle vulnerabilità trattate in
questa sezione.Le dieci categorie che compongono l'edizione 2023 riflettono proprio questa
logica: non genericità, ma rischi specifici del modo in cui un'API espone
dati e funzionalità. Di seguito la lista completa, in ordine di gravità
secondo OWASP\cite{owasp_api_top10_2023}:
\begin{enumerate}
  \item \textbf{API1:2023 -- Broken Object Level Authorization (BOLA)} — un
  endpoint accetta l'ID di una risorsa senza verificare che chi lo richiede
  abbia davvero il diritto di accedervi; basta cambiare un numero nell'URL
  per leggere dati altrui.

  \item \textbf{API2:2023 -- Broken Authentication} — meccanismi di
  autenticazione implementati male, token che non scadono mai, o password
  recuperabili senza controlli sufficienti.

  \item \textbf{API3:2023 -- Broken Object Property Level Authorization} —
  qui il problema non è accedere alla risorsa sbagliata, ma vedere o
  modificare campi che non dovrebbero essere esposti o scrivibili da quel
  particolare utente.

  \item \textbf{API4:2023 -- Unrestricted Resource Consumption} — assenza di
  limiti su chiamate, dimensione dei payload o uso di risorse computazionali,
  con conseguenze che vanno dal semplice rallentamento a un vero e proprio
  denial of service.

  \item \textbf{API5:2023 -- Broken Function Level Authorization} — un
  utente normale riesce a invocare funzioni riservate agli amministratori,
  semplicemente perché il controllo sui permessi manca a livello di endpoint
  o di metodo HTTP.

  \item \textbf{API6:2023 -- Unrestricted Access to Sensitive Business
  Flows} — la singola richiesta può essere anche legittima, ma se ripetuta
  migliaia di volte automatizza un abuso della logica di business (es.
  acquisto massivo di biglietti, creazione automatica di account).

  \item \textbf{API7:2023 -- Server Side Request Forgery (SSRF)} — l'API
  accetta un URL fornito dall'utente e lo interroga senza validarlo,
  permettendo di raggiungere risorse interne altrimenti non accessibili
  dall'esterno.

  \item \textbf{API8:2023 -- Security Misconfiguration} — header di
  sicurezza mancanti, CORS troppo permissivo, ambienti di debug ancora
  attivi in produzione: stessa famiglia di errori vista nel capitolo
  precedente, qui applicata allo strato applicativo.

  \item \textbf{API9:2023 -- Improper Inventory Management} — endpoint
  vecchi, versioni deprecate o ambienti di staging dimenticati online, mai
  rimossi né monitorati, restano comunque raggiungibili da chi li trova.

  \item \textbf{API10:2023 -- Unsafe Consumption of APIs} — la vulnerabilità
  non nasce nell'API stessa, ma nella fiducia cieca riposta nei dati
  restituiti da API di terze parti integrate nel sistema.
\end{enumerate}

Tra le dieci categorie appena elencate, la prima merita un'attenzione
diversa dalle altre. Non tanto per la posizione in classifica, quanto per
quello che rappresenta: un controllo di accesso che manca proprio dove ci
si aspetterebbe fosse più solido, cioè sulla singola risorsa richiesta. È il
tipo di errore che si insinua facilmente in un endpoint scritto in fretta, e
che allo stesso tempo può esporre l'intero set di dati di un sistema con
una richiesta HTTP modificata di poco.

La sezione seguente approfondisce nel dettaglio Broken Object Level
Authorization: come si manifesta, perché resta la vulnerabilità più diffusa
tra quelle riscontrate nelle API reali, e possibili scenari di rimedio.


\subsection{API1:2023 -- Broken Object Level Authorization}

\paragraph{Definizione e inquadramento concettuale}

La Broken Object Level Authorization (BOLA) si manifesta quando un'applicazione, pur autenticando
correttamente l'utente che effettua una richiesta, non verifica se quello stesso utente abbia
effettivamente il diritto di accedere all'oggetto specifico richiesto.

Un'interfaccia di programmazione applicativa espone tipicamente un insieme di risorse --- utenti,
ordini, documenti, transazioni --- identificabili singolarmente attraverso un riferimento univoco.
L'accesso a tali risorse presuppone due verifiche concettualmente distinte, spesso confuse nella
pratica implementativa: l'autenticazione, che accerta l'identità del soggetto che effettua la
richiesta, e l'autorizzazione, che stabilisce se quel soggetto, una volta identificato, abbia il
diritto di operare su una specifica risorsa.

Un utente autenticato correttamente non è, per definizione, un utente autorizzato ad accedere a
qualunque oggetto gestito dal sistema. Questa distinzione, apparentemente elementare, è proprio il
punto in cui molte implementazioni falliscono: il possesso di credenziali valide viene talvolta
considerato --- implicitamente, nel disegno del codice --- condizione sufficiente per l'accesso a
una risorsa, quando in realtà rappresenta soltanto la premessa per una verifica ulteriore, relativa
alla relazione specifica tra quell'utente e quell'oggetto.

\paragraph{Funzionamento delle API e ruolo degli identificatori di risorsa}

Le API REST, e in misura diversa quelle basate su GraphQL, espongono le risorse attraverso
identificatori che compaiono nel percorso della richiesta, nei parametri di query, nel corpo del
messaggio o negli header. Un endpoint come \texttt{GET /api/orders/482} restituisce, nella maggior
parte delle implementazioni, l'oggetto ordine identificato dal valore 482, indipendentemente dal
fatto che tale valore sia un intero progressivo, un UUID o un altro tipo di riferimento.

La gestione di una richiesta di questo tipo comporta, lato server, due passaggi distinti: la
risoluzione dell'identificatore in un oggetto concreto, tramite interrogazione del database o di un
servizio dati, e il controllo che l'utente autenticato abbia diritto di accesso a quell'oggetto
specifico. Il primo passaggio è quasi sempre implementato correttamente, poiché è indispensabile al
funzionamento stesso dell'endpoint; il secondo, non avendo un effetto visibile quando manca, viene
più facilmente trascurato o dato per scontato in fase di sviluppo.

Va osservato che l'endpoint, in questo scenario, non è di per sé vietato all'utente: la sua
configurazione di ruoli e permessi lo autorizza a invocare quella funzione API. Questo distingue
concettualmente BOLA da un'altra categoria della stessa famiglia, la Broken Function Level
Authorization (BFLA), in cui è l'accesso alla funzione stessa a non essere autorizzato. In BOLA la
funzione è legittima; il problema riguarda esclusivamente quale istanza dell'oggetto venga
restituita o modificata.

\paragraph{Meccanismo dell'attacco}

Lo sfruttamento di una vulnerabilità BOLA non richiede, nella maggior parte dei casi, strumenti
sofisticati né la compromissione di credenziali altrui. È sufficiente disporre di un account
legittimo sul sistema bersaglio. Il percorso tipico dell'attacco segue una sequenza riconoscibile:
l'utente accede al proprio account, richiede una propria risorsa attraverso l'interfaccia normale
dell'applicazione e osserva, tramite gli strumenti di sviluppo del browser o un proxy di
intercettazione, l'identificatore associato a quella risorsa nella richiesta sottostante.

A questo punto l'identificatore viene sostituito con un altro valore plausibile per lo stesso
schema --- un intero vicino, nel caso di sequenze numeriche, oppure un valore ottenuto per
enumerazione più sistematica quando lo schema lo consente --- e la richiesta viene reinviata con lo
stesso token di autenticazione, valido e non modificato. Ad esempio:

\begin{verbatim}
GET /api/orders/483 HTTP/1.1
Host: api.example.com
Authorization: Bearer <token_utente_A>
\end{verbatim}

dove 483 non è l'ordine appartenente al titolare del token, ma un identificatore vicino a quello
osservato in precedenza durante il normale utilizzo dell'applicazione. La richiesta è, dal punto di
vista del protocollo, del tutto regolare: cambia soltanto il valore del parametro. Se la risorsa
restituita appartenga realmente a un altro utente dipende interamente da un unico fattore,
approfondito nel paragrafo seguente: se il server esegua o meno un confronto tra l'identità del
chiamante e il proprietario effettivo dell'oggetto.

\paragraph{Esempio applicativo}

Si consideri il caso concreto della risorsa \texttt{/api/orders/\{id\}}, utilizzata per recuperare
i dettagli di un ordine. Una richiesta legittima, effettuata dal proprietario dell'ordine 482, ha
la forma:

\begin{verbatim}
GET /api/orders/482 HTTP/1.1
Host: api.example.com
Authorization: Bearer <token_utente_A>
\end{verbatim}

e produce una risposta simile alla seguente:

\begin{verbatim}
HTTP/1.1 200 OK
Content-Type: application/json

{
  "order_id": 482,
  "customer_id": 1041,
  "items": ["..."],
  "shipping_address": "..."
}
\end{verbatim}

La richiesta modificata introdotta nel paragrafo precedente, diretta all'ordine 483 e appartenente
a \texttt{customer\_id} 1078, viene gestita da un'implementazione vulnerabile allo stesso modo:
l'identificatore viene estratto dal percorso, la tabella corrispondente viene interrogata, il
risultato viene restituito. In nessun punto di questo flusso il valore di \texttt{customer\_id}
restituito viene confrontato con l'identità associata al token presente nell'header
\texttt{Authorization}. La risposta, di conseguenza, è ancora un \texttt{200 OK}, questa volta con
i dati di un ordine che non appartiene al chiamante.

Il punto critico non è la struttura della richiesta, identica nei due casi, ma l'assenza di un
confronto esplicito fra utente autenticato e risorsa richiesta. Un'implementazione corretta
dovrebbe, prima di restituire l'oggetto, verificare che il \texttt{customer\_id} dell'ordine
coincida con l'identità del chiamante --- o che quest'ultimo disponga di un diritto di accesso
esplicitamente concesso su quella risorsa --- e restituire, in caso contrario, un errore di
autorizzazione: tipicamente \texttt{403 Forbidden}, oppure \texttt{404 Not Found} quando si
preferisce non rivelare l'esistenza della risorsa a chi non ne ha diritto.

\paragraph{Cause tecniche ricorrenti}

L'assenza di un controllo di autorizzazione a livello di oggetto è la causa immediata di ogni
istanza di BOLA, ma raramente costituisce una scelta progettuale deliberata. Più spesso deriva da
alcune condizioni che meritano di essere analizzate singolarmente.

La prima riguarda la sovrapposizione concettuale tra autenticazione e autorizzazione nella fase di
implementazione. Un endpoint protetto da un middleware di autenticazione, che verifica la validità
del token e ne estrae l'identità, può apparire sicuro a chi lo sviluppa, inducendo a trascurare la
necessità di un controllo ulteriore specifico per l'oggetto acceduto. Il middleware risponde alla
domanda ``chi sta chiamando?'', non a ``questo chiamante ha diritto su questo oggetto?''.

Una seconda causa è la fiducia implicita nei parametri forniti dal client. Alcune implementazioni
derivano l'identità del proprietario di una risorsa dallo stesso identificatore fornito nella
richiesta, oppure presuppongono che, poiché l'interfaccia utente dell'applicazione non espone mai
identificatori altrui, un client ``onesto'' non li invierebbe mai. Questa assunzione ignora che un
client HTTP può essere costruito o modificato arbitrariamente da chi effettua la richiesta,
indipendentemente dall'interfaccia grafica ufficiale.

Un terzo fattore è rappresentato da controlli di autorizzazione implementati in modo incompleto o
incoerente tra i diversi endpoint di una stessa API. È relativamente comune riscontrare verifiche
corrette sull'operazione di cancellazione di una risorsa e verifiche assenti sull'operazione di
lettura o aggiornamento della medesima risorsa, poiché ciascun endpoint viene sviluppato e
sottoposto a revisione in momenti diversi, con criteri non sempre uniformi. La crescita organica di
un'API nel tempo, con l'aggiunta progressiva di nuovi endpoint da parte di sviluppatori diversi,
tende ad amplificare questa disomogeneità.

Una causa più strutturale riguarda infine la gestione delle relazioni tra utenti e risorse nel
modello dati. Quando tale relazione non è rappresentata in modo esplicito e interrogabile --- ad
esempio tramite una chiave esterna diretta tra la tabella degli utenti e quella delle risorse, o
tramite un livello di autorizzazione centralizzato --- ogni endpoint deve reimplementare
autonomamente la logica di verifica, e questo aumenta la probabilità che qualche percorso venga
trascurato.

\paragraph{Impatto e conseguenze}

L'impatto di una vulnerabilità BOLA dipende dalla natura della risorsa esposta, ma incide
tipicamente su più proprietà di sicurezza contemporaneamente. Sul piano della riservatezza,
l'accesso non autorizzato a un oggetto consente la lettura di dati che dovrebbero rimanere
riservati al proprietario legittimo: informazioni personali, dettagli finanziari, contenuti
sanitari, comunicazioni private, a seconda del dominio applicativo.

Quando l'endpoint vulnerabile consente anche operazioni di scrittura, l'impatto si estende
all'integrità delle risorse: un attaccante può modificare dati appartenenti ad altri utenti,
alterare lo stato di un ordine, cambiare le impostazioni di un account che non gli appartiene. Se
l'operazione esposta è una cancellazione, la conseguenza diventa la perdita di disponibilità della
risorsa per il legittimo proprietario.

Un aspetto che distingue BOLA da molte altre classi di vulnerabilità è la sua natura
sistematicamente ripetibile: una volta individuato un endpoint privo del controllo appropriato, lo
stesso schema di attacco può essere applicato a un numero elevato di identificatori, trasformando
quella che sembra una singola falla puntuale in un canale di accesso a porzioni estese della base
dati. Diversi incidenti di sicurezza documentati pubblicamente, inclusi casi in cui decine di
milioni di record sono stati esposti, hanno questa origine comune.

Il rischio assume connotati particolari negli ambienti multi-tenant e cloud, dove più clienti
condividono la stessa infrastruttura applicativa e la separazione tra i rispettivi dati dipende
interamente dalla correttezza della logica di autorizzazione, non da una separazione fisica delle
risorse. In un simile contesto, una singola vulnerabilità BOLA non compromette soltanto i dati di
un utente, ma può erodere il confine logico tra tenant distinti, con implicazioni che vanno oltre
il singolo incidente e toccano la fiducia complessiva riposta nell'infrastruttura condivisa.

\paragraph{Strategie di mitigazione}

La prevenzione di BOLA richiede che ogni accesso a una risorsa identificata da un parametro
fornito dal client sia accompagnato da una verifica esplicita del diritto dell'utente su quello
specifico oggetto, eseguita interamente lato server. Nessuna informazione proveniente dal client,
incluso l'identificatore stesso della risorsa, può essere considerata prova di autorizzazione:
l'identificatore indica quale oggetto viene richiesto, non chi ha diritto di accedervi.

Un principio applicabile in modo trasversale è quello del least privilege, secondo cui un soggetto
dovrebbe disporre soltanto dei permessi strettamente necessari a svolgere le proprie funzioni,
senza margini di accesso più ampi per comodità implementativa. Applicato al livello oggetto,
questo si traduce nella necessità di associare esplicitamente ogni utente, o ruolo, alle risorse su
cui ha effettivamente diritto di operare, evitando che l'accesso derivi da assunzioni implicite
sulla struttura delle richieste.

Dal punto di vista architetturale, è preferibile centralizzare la logica di autorizzazione in un
livello comune --- un modulo, un middleware o un servizio dedicato --- piuttosto che
reimplementarla in modo indipendente in ciascun endpoint. Questa scelta non elimina la necessità di
applicare il controllo in ogni punto di accesso, ma riduce la probabilità che errori o
dimenticanze locali producano percorsi non protetti, e semplifica la verifica di coerenza tra le
diverse funzioni esposte dall'API.

La coerenza dei controlli tra operazioni di lettura, scrittura e cancellazione sulla stessa risorsa
merita infine attenzione specifica: la protezione di un'operazione non implica automaticamente la
protezione delle altre, ed è proprio in questa asimmetria che si annidano molte delle istanze reali
di BOLA documentate in letteratura.

Nei paragrafi seguenti vengono trattate più sinteticamente le restanti vulnerabilità della OWASP
API Security Top 10, a partire da API2:2023.

\newpage
\subsection{API2:2023 -- Broken Authentication}

\paragraph{Definizione e collocazione}
Broken Authentication occupa la seconda posizione della classificazione OWASP  e si colloca concettualmente a monte di BOLA: mentre quest'ultima riguarda il controllo sull'accesso a uno specifico dato o risorsa, Broken Authentication compromette la capacità stessa del sistema di stabilire con certezza l'identità di chi effettua la richiesta. Un fallimento in questa fase rende vano qualsiasi controllo successivo, poiché un meccanismo di autorizzazione perde ogni efficacia se l'identità su cui si fonda può essere falsificata o aggirata \cite{owasp_api2_2023}.
A differenza di altre categorie più circoscritte, questa voce raggruppa un insieme eterogeneo di debolezze nei processi di verifica dell'identità: dalla gestione poco attenta dei token di sessione all'assenza di vincoli sul numero di tentativi di accesso consentiti. Un esempio comune riguarda i messaggi di risposta del server: se il sistema segnala in modo esplicito se a essere errato sia lo username oppure la password, fornisce involontariamente all'attaccante una conferma sull'esistenza di un dato account. Parallelamente, l'assenza di limitazioni sulla frequenza delle richieste lascia l'interfaccia di login esposta ad attacchi automatizzati volti a indovinare le credenziali su larga scala.
\paragraph{Cause e meccanismi tecnici tipici}
Spesso queste problematiche nascono dal tentativo di implementare soluzioni di autenticazione proprietarie o non standard. Come evidenziato dalla documentazione OWASP, strumenti come le API key o protocolli di autorizzazione come OAuth vengono talvolta impiegati impropriamente per identificare gli utenti in scenari per cui non sono stati progettati, creando falle strutturali \cite{owasp_api2_2023}.
Un ulteriore fattore critico è la mancanza di controlli sul volume dei tentativi di accesso. Endpoint di autenticazione e flussi di recupero password privi di rate limiting risultano vulnerabili ad attacchi in cui elenchi di credenziali già compromesse altrove vengono testati a raffica, sfruttando la tendenza diffusa a riutilizzare le medesime password su più servizi \cite{owasp_credential_stuffing}.
Infine, la gestione dei token di sessione (come i JSON Web Token, o JWT) introduce ulteriori insidie: errori nella convalida della firma crittografica, chiavi di cifratura deboli, mancato controllo dei tempi di scadenza o assenza di meccanismi di revoca immediata permettono a un attaccante di forgiare token validi o di continuare a utilizzare credenziali che dovrebbero essere già decadute \cite{rfc8725}.
\paragraph{Impatto essenziale}
L'impatto di Broken Authentication è per sua natura sistemico: assumendo il controllo dell'identità di un utente, l'attaccante ne eredita tutti i privilegi operativi senza che il backend possa distinguere le azioni legittime da quelle malevole \cite{owasp_api2_2023}. Ciò può tradursi non solo nell'accesso a dati riservati, ma anche nella manomissione dei parametri di sicurezza dell'account (come la modifica dell'email di recupero), trasformando un'intrusione temporanea nella perdita definitiva del profilo.
\paragraph{Mitigazione}
La contromisura principale consiste nell'affidarsi a standard industriali consolidati e a librerie verificate, evitando soluzioni "fai-da-te" per la gestione di token e password \cite{nist_sp800_63b}. A livello applicativo e infrastrutturale è indispensabile introdurre politiche di rate limiting dedicate e restrittive sugli endpoint di login, adottare l'autenticazione a più fattori (MFA) dove possibile e imporre una nuova conferma delle credenziali prima di eseguire modifiche critiche sul profilo utente \cite{owasp_authentication_cheatsheet}.
\paragraph{Rilevabilità}
Sul piano della rilevabilità, l'analisi automatica si divide principalmente in due direttrici a seconda della natura del difetto:
\begin{enumerate}
 \item Analisi dei Token e della Crittografia:
 Per quanto riguarda la gestione dei JWT e dei formati di sessione, le vulnerabilità sono per lo più legate a regole deterministiche e ben definite (come la mancata verifica della firma o l'uso di algoritmi insicuri). La ricerca e gli strumenti di settore riescono quindi a scovare queste falle in modo efficace, sia analizzando staticamente il codice sorgente per intercettare configurazioni errate delle librerie crittografiche, sia generando dinamicamente token manipolati per verificare se il backend li accetti erroneamente \cite{jwteemo, jwtkey}.
\item Flussi di Autenticazione e Abusi Comportamentali: 
Al contrario, individuare l'assenza di rate limiting o la vulnerabilità a tentativi di accesso massivi richiede un approccio basato sul monitoraggio a runtime. Più che da un'ispezione sintattica del codice, queste carenze emergono dall'analisi dinamica del traffico o dall'osservazione dei registri di sistema (log), rilevando picchi anomali di richieste fallite concentrate in brevi finestre temporali \cite{owasp_credential_stuffing}.
\end{enumerate}







\subsection{API3:2023 -- Broken Object Property Level Authorization}

\paragraph{Definizione e collocazione}

La \textit{Broken Object Property Level Authorization} (BOPLA) identifica una classe di vulnerabilità nella quale i controlli di autorizzazione applicati alle API non considerano adeguatamente le singole proprietà degli oggetti gestiti dall'applicazione. Il problema non riguarda quindi soltanto la possibilità per un utente di accedere a un determinato oggetto, ma anche quali attributi di quell'oggetto possa effettivamente leggere o modificare.


La distinzione rispetto alla \textit{Broken Object Level Authorization} (BOLA) è rilevante. Nel caso di BOLA, il controllo di autorizzazione fallisce a livello dell'oggetto: un utente riesce, per esempio, ad accedere a una risorsa appartenente a un altro utente manipolando l'identificativo della risorsa. Nel caso di BOPLA, invece, l'utente può avere legittimamente accesso all'oggetto, ma ottenere o modificare proprietà che non rientrano nei suoi privilegi. L'oggetto, quindi, non costituisce necessariamente la risorsa a cui viene negato l'accesso; è una sua specifica proprietà a essere esposta o manipolabile senza un'adeguata autorizzazione.

Il rischio è particolarmente rilevante nelle API REST, dove è frequente che un endpoint restituisca direttamente la rappresentazione serializzata di un oggetto. Una struttura interna contenente numerosi attributi può così diventare, senza una selezione esplicita, la struttura della risposta esposta al client.

\paragraph{Cause e meccanismi tecnici tipici}

Una delle cause più comuni è la sovrapposizione tra il modello utilizzato internamente dall'applicazione e quello esposto attraverso l'API. Un oggetto memorizzato nel database può contenere, ad esempio, proprietà quali \texttt{id}, \texttt{name}, \texttt{email}, \texttt{role} e \texttt{internalStatus}. Non tutte queste proprietà hanno però la stessa finalità né devono necessariamente essere accessibili allo stesso insieme di utenti.

Un'implementazione che serializza automaticamente l'intero oggetto può produrre una risposta contenente anche attributi che il client non dovrebbe conoscere. In questo caso il problema non deriva necessariamente da un accesso non autorizzato all'oggetto, ma dal fatto che l'API non applica una politica distinta alle proprietà che lo compongono. L'interfaccia grafica potrebbe nascondere determinati campi, ma ciò non costituisce una misura di sicurezza: il contenuto della risposta HTTP rimane osservabile dal client e può essere analizzato direttamente. 

Un secondo meccanismo riguarda le operazioni di scrittura. Framework e librerie possono consentire il \textit{binding} automatico dei parametri ricevuti dal client verso gli attributi di un oggetto applicativo. Se il server accetta un payload come:

\begin{verbatim}
{
"name": "Mario",
"email": "[mario@example.com](mailto:mario@example.com)"
}
\end{verbatim}

un attaccante può provare ad aggiungere proprietà che non erano previste dal normale flusso applicativo:

\begin{verbatim}
{
"name": "Mario",
"email": "[mario@example.com](mailto:mario@example.com)",
"role": "admin"
}
\end{verbatim}

Se il campo \texttt{role} viene accettato e aggiornato senza un controllo specifico, l'utente può modificare una proprietà interna che dovrebbe essere gestita esclusivamente dal sistema. Lo stesso principio può riguardare attributi quali \texttt{isAdmin}, \texttt{approved}, \texttt{price} o \texttt{ownerId}, a seconda del modello utilizzato dall'applicazione. 

Il problema può quindi manifestarsi in due direzioni differenti: lettura non autorizzata di proprietà e modifica non autorizzata di proprietà. Nel primo caso l'attaccante osserva informazioni che non dovrebbe ricevere; nel secondo riesce a influenzare direttamente lo stato dell'applicazione attraverso attributi che non dovrebbero essere controllabili dal client.

\paragraph{Impatto essenziale}

L'impatto di una vulnerabilità BOPLA dipende principalmente dal tipo di proprietà coinvolta e dalle funzioni che la utilizzano. La semplice esposizione di un attributo può comportare una violazione della riservatezza, mentre la possibilità di modificarlo può compromettere l'integrità dei dati o il funzionamento dell'applicazione.

Nel caso di proprietà contenenti informazioni personali, finanziarie o relative all'organizzazione interna del sistema, l'esposizione può determinare una divulgazione di dati a soggetti non autorizzati. La possibilità di modificare determinate proprietà può invece portare alla perdita o alla corruzione dei dati. In situazioni più critiche, una proprietà utilizzata nei meccanismi di autorizzazione o nella gestione dello stato dell'account può essere sfruttata per ottenere un'escalation dei privilegi o contribuire alla compromissione dell'account.


\textbf{La gravità della vulnerabilità dipende quindi dal significato assunto dalla proprietà all'interno dell'applicazione
}

\paragraph{Mitigazione}

La mitigazione di BOPLA richiede che l'autorizzazione venga applicata non soltanto all'endpoint o all'oggetto, ma anche alle proprietà che possono essere restituite o modificate. La prima misura consiste quindi nel definire esplicitamente quali attributi ogni endpoint può esporre e quali possono essere accettati in input.

Per le risposte, è preferibile evitare l'utilizzo di serializzazione quando questi comportano la restituzione automatica dell'intero oggetto. È invece opportuno selezionare esplicitamente le proprietà necessarie alla specifica operazione.

Per le operazioni di scrittura è necessario limitare esplicitamente gli attributi modificabili dal client. Una tecnica comune consiste nell'utilizzare una \textit{whitelist} delle proprietà ammesse, evitando che il payload venga associato automaticamente a tutti gli attributi dell'oggetto. 


\paragraph{Rilevabilità}
La rilevazione di BOPLA si articola in due verifiche distinte, corrispondenti alle due direzioni in cui una proprietà può attraversare il confine fra client e server.

La prima riguarda le proprietà esposte in risposta a una richiesta legittima, e si riduce al confronto fra gli attributi effettivamente restituiti e quelli previsti per il ruolo o il contesto dell'utente. È la verifica più diretta, perché l'evidenza è contenuta nella risposta stessa.

La seconda riguarda le proprietà accettate in ingresso, ed è più onerosa perché il server non le dichiara. Individuarle richiede di modificare il payload delle richieste, eventualmente ricorrendo a tecniche di fuzzing per scoprire attributi non documentati, e di verificare poi se il server ne recepisca effettivamente il valore.

Le due verifiche non sono sostituibili l'una con l'altra: un'API può accettare in input una proprietà sensibile che non restituisce mai, oppure restituire dati riservati senza consentirne la modifica. Un'analisi che si limitasse a una sola direzione lascerebbe scoperto metà del problema.


\subsection{API4:2023 -- Unrestricted Resource Consumption}
\label{subsubsec:api4}

\paragraph{Definizione e collocazione}
Ogni richiesta servita da un'API consuma risorse: banda, CPU, memoria, storage e, con frequenza crescente,la categoria API4:2023 raccoglie i casi in cui almeno uno dei limiti che dovrebbero governare tale consumo è assente o dimensionato in modo inadeguato~\cite{owasp_api4_2023}. Il documento OWASP elenca i limiti attesi: timeout di esecuzione, memoria allocabile, numero di descrittori di file e di processi,insieme a molte altre categorie.~\cite{owasp_api4_2023}.


La collocazione nella Top~10 richiede un chiarimento ulteriore, perché il consumo di risorse non è un difetto di una specifica funzione ma una proprietà trasversale di qualunque endpoint, e i limiti che lo governano sono la difesa in profondità di altre categorie. Quattro relazioni contano più delle altre.

Con API2:2023 il rapporto è di specializzazione. Broken Authentication chiede
che gli endpoint di autenticazione e di recupero credenziali siano protetti da
meccanismi anti-brute-force più severi del rate limiting ordinario applicato al
resto dell'API~\cite{owasp_api2_2023}. Il meccanismo di aggiramento è però lo
stesso che API4 elenca fra i limiti mancanti, ossia il numero di operazioni
ammesse in una singola richiesta come nel batching
GraphQL~\cite{owasp_api4_2023}. A cambiare è il bene protetto, e i due scenari
illustrativi lo mostrano bene: in API2 il batch aggrega tentativi di login per
superare un limite di tre richieste al minuto e arrivare alla
credenziale~\cite{owasp_api2_2023}, in API4 aggrega caricamenti di immagini fino
a esaurire la memoria del server~\cite{owasp_api4_2023}.

Con API6:2023 il vettore resta l'automazione delle richieste e cambia l'esito.
OWASP colloca il danno sul piano del business e scrive, nella descrizione degli
impatti, che un impatto tecnico in generale non è atteso~\cite{owasp_api6_2023}.
Lo scenario di apertura descrive l'acquisto automatico dell'intero stock di una
console distribuito su indirizzi IP e località
diverse~\cite{owasp_api6_2023}: una costruzione che aggira per definizione il
throttling per IP, cioè la forma di limitazione più semplice da implementare e
anche la più fragile.

Con API10:2023 la sovrapposizione è nel testo stesso. Fra gli indicatori di
Unsafe Consumption of APIs compaiono l'assenza di un limite alle risorse
impiegate per elaborare le risposte dei servizi terzi e l'assenza di timeout
nelle interazioni con essi~\cite{owasp_api10_2023}, due condizioni che API4
annovera già fra i limiti la cui mancanza rende un'API
vulnerabile~\cite{owasp_api4_2023}. A distinguerle è la direzione del consumo:
API4 guarda alle risorse che un client sottrae al server, API10 a quelle che un
fornitore a monte può far spendere al server.

Con API9:2023 il problema non è l'assenza del limite ma la sua copertura
parziale. Lo scenario di Improper Inventory Management descrive un rate
limiting collocato non nel codice dell'API ma in un componente separato fra
client e API ufficiale, e un host beta che espone la stessa API senza quel
componente: da lì è stato possibile forzare per tentativi il token di reset a
sei cifre e cambiare la password di qualunque utente~\cite{owasp_api9_2023}.
OWASP ne ricava due indicazioni operative, documentare il rate limiting fra gli
aspetti del contratto API ed estendere le misure di protezione esterne a tutte
le versioni esposte e non solo a quella in produzione~\cite{owasp_api9_2023}.

\begin{table}[H]
\centering
\small
\begin{tabular}{p{0.11\textwidth} p{0.38\textwidth} p{0.38\textwidth}}
\hline
\textbf{Categoria} & \textbf{Punto di contatto con API4} & \textbf{Elemento distintivo} \\
\hline
API2 & Gli endpoint di autenticazione e recupero credenziali richiedono limiti più severi di quelli ordinari; il batching GraphQL aggira il rate limiting in entrambe le categorie & Bene protetto: la segretezza della credenziale (CWE-307), dove API4 tutela disponibilità e costo (CWE-770, CWE-799) \\
API6 & Automazione di richieste di per sé legittime; distribuendo l'attacco si elude il throttling per indirizzo IP & Danno collocato sul piano del business, con impatto tecnico in generale non atteso; contromisure comportamentali \\
API9 & Il limite può risiedere in un componente esterno al codice dell'API e non coprire tutti gli host esposti & Difetto di inventario e documentazione: il codice resta corretto mentre la superficie esposta rimane scoperta \\
API10 & Assenza di timeout e di limiti sulle risorse impiegate per elaborare le risposte di servizi terzi & Direzione del consumo: le risorse che un fornitore a monte fa spendere al server \\
\hline
\end{tabular}
\caption{Relazioni fra API4:2023 e le categorie contigue della OWASP API Security Top~10 2023.}
\label{tab:api4_relazioni}
\end{table}

\paragraph{Cause e meccanismi tecnici tipici}
La radice comune è un'asimmetria: produrre una richiesta costa al client una quantità trascurabile di risorse, mentre servirla può costare al server una quantità che cresce in funzione di parametri controllati dal client stesso. Ogni volta che questa funzione di costo non è limitata superiormente, l'API è esposta. Conviene distinguere quattro meccanismi, perché ciascuno richiede una difesa diversa e presenta un profilo di rilevabilità distinto, e poi chiedersi perché, in pratica, i limiti vengano omessi.

Il primo meccanismo è il volume puro: l'assenza di un limite alla frequenza delle richieste per client, per utenza o per endpoint. OWASP osserva che lo sfruttamento non richiede competenze particolari, poiché richieste concorrenti possono essere generate da una singola macchina o da risorse cloud noleggiate, e che la maggior parte degli strumenti automatici disponibili mira proprio a saturare il tasso di servizio~\cite{owasp_api4_2023}. La forma più insidiosa non è però l'assenza totale del limite, ma la sua presenza in una sola dimensione: un limite per indirizzo IP è inefficace contro un attacco distribuito, un limite per endpoint è inefficace contro il batching, un limite per utente autenticato non copre gli endpoint anonimi, che sono tipicamente quelli di registrazione e recupero credenziali.

Il secondo è il parametro che governa la dimensione della risposta. Un endpoint paginato che accetta \texttt{limit=100000}, o che non impone alcun tetto, trasforma una singola richiesta in una scansione completa della tabella sottostante. Non a caso OWASP colloca fra le contromisure la validazione lato server dei parametri che controllano il numero di record restituiti~\cite{owasp_api4_2023}. Il caso si intreccia con API3:2023, già trattato: la stessa disattenzione verso la forma della richiesta produce lì esposizione di proprietà, qui esplosione del carico, e un endpoint che soffre di entrambe consente a un attaccante di estrarre l'intero insieme di dati sensibili con una manciata di chiamate.

Il terzo meccanismo è la moltiplicazione di operazioni in una sola richiesta. Lo scenario GraphQL riportato da OWASP è istruttivo: un'API applica rate limiting all'endpoint \texttt{/graphql} e controlla la dimensione delle immagini caricate, ma non limita il numero di mutazioni aggregate in un unico batch; l'attaccante invia centinaia di \texttt{uploadPic} in una richiesta sola, ciascuna delle quali innesca la generazione di miniature, e satura la memoria del server aggirando entrambe le protezioni~\cite{owasp_api4_2023}. La letteratura ha formalizzato il fenomeno: Cha et~al. mostrano che una query GraphQL nel caso peggiore impone al server una quantità di lavoro esponenziale e ricordano, richiamando lo studio empirico precedente dello stesso gruppo, che la maggioranza degli schemi pubblici espone il fornitore a query ad alto costo~\cite{cha2020graphql}. Il meccanismo non è esclusivo di GraphQL: qualunque endpoint REST che accetti un array di operazioni, come un'importazione massiva o un aggiornamento in blocco, presenta lo stesso profilo se non limita la cardinalità dell'array.

Il quarto meccanismo non dipende dal numero di richieste ma dalla complessità algoritmica dell'elaborazione di un singolo input. È la famiglia degli attacchi di complessità algoritmica descritta da Crosby e Wallach~\cite{crosby2003algorithmic}, di cui il ReDoS (\emph{Regular expression Denial of Service}) è l'istanza più diffusa nel contesto delle API: un'espressione regolare soggetta a backtracking catastrofico, applicata a un input costruito ad arte, può bloccare un thread per minuti. Staicu e Pradel hanno documentato la gravità del fenomeno nei server basati su JavaScript, la cui esecuzione a thread singolo rende un'unica richiesta sufficiente a congelare l'intero processo~\cite{staicu2018redos}. Nella stessa famiglia rientrano i timeout mancanti sulle chiamate in uscita verso altri servizi, che consentono a un fornitore lento di trattenere indefinitamente connessioni e worker.

A questi si aggiunge la causa che l'edizione 2023 porta in primo piano: l'assenza di soglie di spesa verso i fornitori esterni. Qui la risorsa consumata non è computazionale ma economica, e il limite mancante non risiede nel codice ma nel contratto con il provider~\cite{owasp_api4_2023}.

Resta da spiegare perché limiti così semplici da enunciare vengano omessi con la frequenza che OWASP classifica come diffusa. Tre ragioni ricorrono. La prima è la diffusione della responsabilità: un limite può essere applicato dal framework, dal reverse proxy, dal gateway, dall'orchestratore o dal fornitore cloud, e quando ogni livello presume che lo faccia un altro il risultato è che nessuno lo fa. La seconda è che il difetto non produce errori: un endpoint senza limiti funziona perfettamente sotto carico normale e supera ogni test funzionale, sicché nulla nel ciclo di sviluppo lo segnala. La terza è l'elasticità dell'infrastruttura, che nasconde il sintomo trasformando l'esaurimento in scalabilità automatica e quindi in costo; il problema diventa visibile alla fattura successiva, non nel monitoraggio.

\paragraph{Impatto essenziale}
L'esito più immediato è il denial of service per esaurimento. CWE-770 lo enumera come conseguenza principale sulla disponibilità, per CPU, memoria o altre risorse, notando che l'attaccante può ottenerlo sia con molte richieste rapide sia inducendo l'uso di risorse più grandi del necessario~\cite{cwe770}. Diversamente da un attacco volumetrico di rete, il punto di rottura non è la banda ma la logica applicativa, e la quantità di traffico necessaria può essere di ordini di grandezza inferiore: una singola richiesta ReDoS o un solo batch GraphQL sono sufficienti, e nessuna difesa perimetrale basata sul volume li riconosce come anomali.

L'edizione 2023 insiste però su una seconda dimensione, quella economica, e le riconosce pari dignità. Due dei tre scenari illustrativi non causano alcuna interruzione del servizio. Nel primo, un flusso di recupero password che invia un codice via SMS attraverso un fornitore che tariffa 0,05 dollari a messaggio viene invocato decine di migliaia di volte da uno script, con perdite di migliaia di dollari in pochi minuti. Nel terzo, un file distribuito tramite API supera la soglia di 15\,GB oltre la quale la cache non opera e, in assenza di allarmi sui costi, la fattura mensile per lo storage passa da circa 13 a circa 8\,000 dollari~\cite{owasp_api4_2023}. Il tratto comune è che l'infrastruttura elastica converte l'esaurimento in spesa: il servizio resta disponibile, ma a un costo che nessuno ha autorizzato. OWASP nota anche che l'attaccante, pur privo di visibilità diretta sui costi, può inferirli dal modello di prezzo del provider~\cite{owasp_api4_2023}.

Un terzo effetto, meno citato, è la degradazione progressiva: prima del collasso, la latenza cresce per tutti gli utenti legittimi e, nelle architetture a microservizi, l'esaurimento di un componente si propaga a monte attraverso code e tentativi ripetuti. In questo senso l'impatto di API4 è raramente confinato all'endpoint vulnerabile: un pool di connessioni al database saturato da una query di paginazione senza tetto rende indisponibili anche gli endpoint correttamente limitati che condividono lo stesso pool.

\paragraph{Mitigazione}
Le contromisure si dispongono su livelli distinti e nessuno, da solo, è sufficiente. Al livello più basso, OWASP raccomanda di affidare i limiti fisici (memoria, CPU, numero di processi e descrittori, riavvii) a piattaforme che li rendono dichiarativi, come container e funzioni serverless~\cite{owasp_api4_2023}. Si tratta di una difesa di ultima istanza: non impedisce l'abuso, lo confina.

Al livello dell'ingresso dei dati, la regola è dichiarare e applicare una dimensione massima per ogni parametro e payload: lunghezza delle stringhe, cardinalità degli array, dimensione dei file e, in particolare, il valore dei parametri che controllano il numero di record restituiti~\cite{owasp_api4_2023}. CWE-770 inquadra la stessa misura come validazione dell'input per lista di valori ammessi, con l'avvertenza che copre solo i casi in cui l'input influenza dimensione o frequenza delle allocazioni~\cite{cwe770}. Il contratto OpenAPI è il luogo naturale per rendere questi vincoli espliciti attraverso \texttt{maxLength}, \texttt{maxItems} e \texttt{maximum} sui parametri di paginazione. Un contratto che dichiara i limiti ha un valore che va oltre la documentazione: rende i limiti verificabili da uno strumento, tema su cui si tornerà.

Al livello delle interazioni, il rate limiting resta la contromisura di riferimento, con due precisazioni ricorrenti nel documento OWASP: va tarato per endpoint in base alle esigenze di business, e va affiancato da un limite sul numero di esecuzioni di una singola operazione per utente (validazione di un OTP, richiesta di reset), che il solo conteggio delle chiamate all'endpoint non cattura~\cite{owasp_api4_2023}. La seconda precisazione è quella che lo scenario GraphQL rende evidente. Per gli endpoint di autenticazione il requisito si irrigidisce: OWASP chiede meccanismi anti-brute-force più severi del rate limiting ordinario e il trattamento degli endpoint di recupero credenziali alla stregua di quelli di login~\cite{owasp_api2_2023}, e le linee guida NIST sull'autenticazione digitale prescrivono al verificatore controlli di throttling contro gli attacchi di indovinamento in linea~\cite{nist_sp800_63b}. Sul piano protocollare, HTTP fornisce il codice di stato \texttt{429 Too Many Requests}, definito in RFC~6585~\cite{rfc6585}, e l'IETF sta standardizzando i campi \texttt{RateLimit} e \texttt{RateLimit-Policy} con cui il server comunica al client quota residua e finestra temporale, consentendo una regolazione cooperativa del traffico~\cite{ietf_ratelimit_headers}; al momento della stesura il documento è ancora una Internet-Draft. NIST tratta il throttling tra le strategie di sicurezza per le applicazioni a microservizi, collocandolo nel gateway o nella mesh piuttosto che nel singolo servizio~\cite{nist_sp800_204}. Quest'ultima indicazione va letta insieme allo scenario dell'host beta già richiamato: collocare il limite nel gateway è corretto purché ogni host che espone l'API stia dietro quel gateway, e il limite applicato nel codice, per quanto meno elegante, ha il vantaggio di viaggiare con l'applicazione in ogni ambiente in cui viene distribuita.

Per la famiglia algoritmica la mitigazione è specifica: timeout su ogni operazione potenzialmente lunga (chiamate in uscita, elaborazione di immagini, matching di espressioni regolari) e, per il ReDoS, l'adozione di motori regex a tempo lineare o la riscrittura delle espressioni vulnerabili~\cite{staicu2018redos}. Per GraphQL, l'analisi statica del costo di una query prima della sua esecuzione, come quella proposta da Cha et~al., permette di rifiutare richieste il cui limite superiore di complessità eccede una soglia~\cite{cha2020graphql}.

Per la dimensione economica, infine, OWASP prescrive soglie di spesa su tutti i fornitori integrati e, dove il provider non le supporta, almeno allarmi di fatturazione~\cite{owasp_api4_2023}. È una misura amministrativa più che tecnica, ma è l'unica che intercetta lo scenario del file da 18\,GB.

CWE-770 aggiunge una considerazione di progetto che vale come avvertenza: il throttling su base identitaria può a sua volta diventare un vettore, perché un attaccante che impersona un utente legittimo può farlo bloccare; il throttling uniforme evita questo problema ma non risolve l'attacco, si limita a renderlo più costoso~\cite{cwe770}. Il compromesso fra le due strategie va deciso caso per caso, e la risposta cambia fra un endpoint di login, dove il blocco dell'account è un rischio concreto di negazione del servizio mirata, e un endpoint di ricerca, dove il throttling per utente non danneggia terzi.

\paragraph{Rilevabilità}
Le due fonti primarie sembrano contraddirsi. OWASP classifica la rilevabilità come facile, motivando che richieste costruite con parametri che controllano il numero di risorse restituite, seguite dall'analisi di stato, tempo e lunghezza della risposta, sono normalmente sufficienti a identificare il problema, e che lo stesso vale per le operazioni batch~\cite{owasp_api4_2023}. CWE-770, nella sezione sui metodi di rilevamento, giudica invece l'analisi statica automatica di utilità limitata per le allocazioni senza limite, con eccezioni circoscritte (mancato rilascio di risorse di sistema, argomenti non verificati alle allocazioni di memoria, file di configurazione privi di un valore massimo), e la ritiene inadatta per le risorse definite dall'applicazione, come una politica che consenta a un utente un numero massimo di operazioni al giorno~\cite{cwe770}. La contraddizione è apparente: OWASP descrive la prospettiva del tester in black box, CWE quella dello strumento che analizza il codice. Le due prospettive vanno trattate separatamente.

Sul versante dinamico l'evidenza è consolidata. CWE-770 riconosce che le tecniche automatiche basate sulla generazione di un numero elevato di richieste in un intervallo breve sono efficaci nel produrre gli effetti collaterali dell'allocazione incontrollata, pur richiedendo interpretazione manuale dei risultati, e che il fuzzing può rilevare la debolezza in modo opportunistico, specialmente quando l'allocazione dipende da input numerici~\cite{cwe770}. Il lavoro di Antunes, Neves e Veríssimo, citato dalla stessa voce CWE, è fra i primi a proporre un approccio sistematico di rilevamento e predizione delle vulnerabilità da esaurimento nei server di rete tramite iniezione di carico~\cite{antunes2008resource}. La metodologia OWASP, ridotta all'essenziale, è una forma di test differenziale: si confrontano le risposte a richieste con parametri di dimensione crescente, e un tempo o una lunghezza che crescono senza saturare rivelano l'assenza del limite. Il test dinamico ha un vantaggio decisivo rispetto a ogni analisi del codice: osserva il sistema nel suo insieme, e quindi vede il limite ovunque sia applicato. Ha però un limite simmetrico: osserva solo l'host che interroga, e un host beta senza gateway resta invisibile finché non viene inventariato.

Sul versante statico la letteratura si concentra su sottoclassi ben delimitate. Per il ReDoS esiste un filone maturo: Staicu e Pradel combinano identificazione statica di espressioni regolari vulnerabili e verifica dinamica su siti reali, mostrando che il problema è diffuso e sfruttabile in produzione~\cite{staicu2018redos}. Per GraphQL, l'analisi di Cha et~al. calcola in tempo lineare un limite superiore del costo di una query senza eseguirla, ed è stata validata su API commerciali confrontando la complessità predetta con quella osservata nelle risposte~\cite{cha2020graphql}; è a tutti gli effetti un rilevatore statico, sebbene operi sulla query in arrivo e non sul codice del server.

Ciò che manca, per quanto è stato possibile verificare, è un corpo di lavori sottoposti a revisione paritaria dedicato al rilevamento statico dell'assenza di rate limiting, di timeout o di limiti di paginazione nel codice sorgente di API REST. Le pratiche esistenti in questo spazio provengono dagli strumenti e non dalla letteratura, e si lasciano raggruppare in tre famiglie di segnali. La prima opera sul codice: regole di pattern matching che segnalano chiamate HTTP in uscita prive di timeout, applicazioni web che non registrano alcun middleware di limitazione, parametri di paginazione letti dalla richiesta e passati alla query senza un tetto, espressioni regolari con struttura vulnerabile applicate a input esterni. La seconda opera sui manifest infrastrutturali: presenza di limiti di risorse per i container, dimensione massima del corpo della richiesta nel reverse proxy, piani di throttling sui gateway. La terza opera sul contratto OpenAPI: schemi privi di \texttt{maxLength}, \texttt{maxItems} o \texttt{maximum}, endpoint che dichiarano parametri di paginazione senza vincoli, assenza di risposte \texttt{429} documentate. L'osservazione di CWE-770 sui file di configurazione privi di un massimo è, di fatto, la sola indicazione proveniente da una fonte primaria che legittima la seconda e la terza direzione~\cite{cwe770}.

Il quadro suggerisce una conclusione precisa sulla natura del segnale statico, che distingue API4 dalle categorie trattate finora. Per una BOLA o un difetto di autenticazione, l'assenza del controllo nel codice del gestore è la vulnerabilità, perché il controllo non può che stare lì. Per API4 l'analisi del codice non può affermare la presenza della vulnerabilità, perché il limite potrebbe essere applicato altrove: nel gateway, nel reverse proxy, nella piattaforma di orchestrazione. Può però affermare con buona confidenza l'assenza di un controllo in un dato livello, e questa assenza diventa un finding solo quando è confermata sugli altri livelli o da una sonda dinamica. Lo scenario di API9 mostra anche il caso opposto: la presenza del limite in un livello non garantisce la copertura di tutti gli host~\cite{owasp_api9_2023}. È il caso in cui la correlazione fra segnali eterogenei, provenienti dal codice, dall'infrastruttura e dal contratto, non è un raffinamento ma una condizione necessaria per formulare un giudizio, ed è il criterio che la parte sperimentale di questo lavoro adotta per la categoria.




\subsection{API5:2023 -- Broken Function Level Authorization}
\label{subsubsec:api5}

\paragraph{Definizione e collocazione}
Un'API espone operazioni, e non tutte sono destinate a tutti gli utenti.
Creare un invito con privilegi amministrativi, esportare l'elenco completo
degli utenti, cancellare un account altrui: sono funzioni che l'applicazione
riserva a un ruolo, e il controllo che verifica il ruolo del chiamante prima di
eseguirle è l'autorizzazione a livello di funzione. API5:2023 raccoglie i casi
in cui tale controllo manca o è aggirabile, e in cui un utente anonimo o
ordinario riesce di conseguenza a invocare funzioni
riservate~\cite{owasp_api5_2023}.

La differenza rispetto ad API1:2023 è di livello, ed è l'unica che OWASP fissi
in modo prescrittivo. Nella pagina dedicata alla BOLA si legge che l'accesso
dell'utente all'endpoint vulnerabile è previsto dal progetto e che la
violazione avviene manipolando l'identificatore dell'oggetto, mentre se
l'attaccante raggiunge un endpoint o una funzione cui non dovrebbe accedere il
caso ricade sotto la BFLA~\cite{owasp_api1_2023}. La domanda da porsi davanti a
un finding non riguarda quindi i dati esposti, ma la legittimità della chiamata
in sé: chi sostituisce un identificatore su un endpoint cui ha legittimamente
accesso e legge il profilo altrui commette una BOLA; chi invoca l'esportazione
completa degli utenti commette una BFLA, anche se il risultato include il
proprio profilo. La letteratura sul rilevamento nomina le due forme come
escalation orizzontale e verticale, intendendo per verticale il caso in cui un
utente ordinario assume privilegi amministrativi e per orizzontale quello in
cui assume i privilegi di un altro utente
ordinario~\cite{monshizadeh2014mace}: API5 è la forma verticale.

Il riferimento tassonomico indicato da OWASP è CWE-285 \emph{Improper
Authorization}~\cite{owasp_api5_2023}, che descrive il prodotto che non esegue,
o esegue scorrettamente, un controllo di autorizzazione quando un attore tenta
di accedere a una risorsa o compiere un'azione~\cite{cwe285}. MITRE ne
sconsiglia l'uso per mappare vulnerabilità reali, trattandosi di una classe di
livello elevato, e rimanda alle figlie: CWE-862 \emph{Missing Authorization}
quando il controllo è assente, CWE-863 \emph{Incorrect Authorization} quando
esiste ma decide male~\cite{cwe285}. La stessa voce dedica una nota
terminologica alla distinzione fra autorizzazione e autenticazione,
raccomandando le abbreviazioni AuthZ e AuthN separate ed evitando la forma
ambigua \emph{Auth}~\cite{cwe285}. La distinzione vale anche qui: un'API può
autenticare in modo impeccabile e restare priva di autorizzazione funzionale, e
un token contraffatto fornisce un'identità, ma è il controllo funzionale
mancante a trasformarla in privilegio.

Oltre ad API1, la categoria tocca altre quattro voci dell'elenco, riassunte
nella Tabella~\ref{tab:api5_relazioni}.

\begin{table}[H]
\centering
\small
\begin{tabular}{p{0.11\textwidth} p{0.38\textwidth} p{0.38\textwidth}}
\hline
\textbf{Categoria} & \textbf{Punto di contatto con API5} & \textbf{Criterio di distinzione} \\
\hline
API1 & Difetti di autorizzazione su un chiamante autenticato & L'accesso all'endpoint è previsto (API1) oppure indebito (API5); demarcazione fissata da OWASP \\
API2 & Precede API5 nella catena: senza identità non c'è ruolo da verificare & Accertamento dell'identità (AuthN) contro verifica dei privilegi (AuthZ) \\
API3 & Il campo che codifica il ruolo nel corpo della richiesta appartiene a entrambe & Scrivibilità della proprietà (API3) contro raggiungibilità dell'operazione (API5) \\
API8 & Verbi HTTP non necessari e discrepanze fra i server della catena aggirano il controllo & Difetto di configurazione dello stack anziché di logica applicativa \\
API9 & Versioni vecchie ed endpoint deprecati espongono funzioni amministrative protette altrove & Difetto di inventario: il controllo esiste ma non copre tutta la superficie \\
\hline
\end{tabular}
\caption{Relazioni fra API5:2023 e le altre categorie della OWASP API Security Top~10 2023.}
\label{tab:api5_relazioni}
\end{table}

\paragraph{Cause e meccanismi tecnici tipici}

La radice è un'asimmetria di costo: produrre una richiesta costa al client
quasi nulla, servirla può costare al server una quantità che cresce con
parametri che il client controlla. Finché quella funzione non ha un limite
superiore l'API resta esposta.

Il caso più semplice è l'assenza di un tetto alla frequenza, che OWASP giudica
sfruttabile senza competenze particolari~\cite{owasp2023api4}. Più insidiosa è
la presenza del limite in una sola dimensione: un tetto per indirizzo IP cade
contro un attacco distribuito, uno per endpoint non intercetta il batching, uno
per utente autenticato lascia scoperti gli endpoint anonimi di registrazione e
recupero credenziali. Altrove il difetto sta nei parametri che governano la
dimensione della risposta, e un endpoint paginato che accetta
\texttt{limit=100000} trasforma una chiamata in una scansione completa della
tabella. Il batching moltiplica invece le operazioni dentro una richiesta sola,
come nello scenario GraphQL in cui centinaia di \texttt{uploadPic} aggregate
saturano la memoria aggirando sia il rate limiting sia il controllo sulle
dimensioni~\cite{owasp2023api4}; Cha et al. hanno mostrato che nel caso
peggiore una query impone al server un lavoro
esponenziale~\cite{cha2020graphql}. Resta la famiglia che dipende dal costo
algoritmico di una singola richiesta~\cite{crosby2003}, di cui il ReDoS è
l'istanza più diffusa nei server a thread singolo~\cite{staicu2018}.

Che limiti tanto semplici manchino dipende da due fatti. Il limite può
risiedere nel framework, nel proxy, nel gateway o presso il fornitore cloud, e
quando ogni livello presume che se ne occupi un altro non se ne occupa nessuno.
Il difetto, poi, non produce errori: un endpoint senza limiti supera ogni test
funzionale, e l'elasticità dell'infrastruttura converte l'esaurimento in costo,
facendo comparire il sintomo sulla fattura anziché nel monitoraggio.

\paragraph{Impatto essenziale}

L'esito immediato è il denial of service per esaurimento, che CWE-770 enumera
come conseguenza principale sulla disponibilità~\cite{cwe770}. Il punto di
rottura è la logica applicativa anziché la banda, e il traffico richiesto può
essere di ordini di grandezza inferiore: una richiesta ReDoS o un solo batch
GraphQL bastano, e nessuna difesa perimetrale basata sul volume li riconosce
come anomali. L'edizione 2023 dà pari dignità alla dimensione economica, e due
dei tre scenari non producono alcuna interruzione: un flusso di recupero
password che invia SMS a 0,05 dollari l'uno, invocato decine di migliaia di
volte, e un file che supera la soglia di cache, con la fattura mensile che
passa da circa 13 a circa 8\,000 dollari~\cite{owasp2023api4}.

\paragraph{Mitigazione}

Nessun livello basta da solo. I limiti fisici su memoria, CPU e processi OWASP
li affida a piattaforme che li rendono dichiarativi~\cite{owasp2023api4},
difesa di ultima istanza che confina l'abuso senza impedirlo. All'ingresso dei
dati la regola è dichiarare una dimensione massima per ogni parametro e
payload, con attenzione a quelli che governano il numero di record
restituiti~\cite{owasp2023api4}; il contratto OpenAPI è il luogo naturale per
esplicitare questi vincoli, che così diventano verificabili da uno strumento.
Sul piano delle interazioni il rate limiting va tarato per endpoint e
affiancato da un limite sulle esecuzioni della singola operazione per utente,
che il conteggio delle chiamate non cattura~\cite{owasp2023api4}, e si
irrigidisce sugli endpoint di autenticazione~\cite{owasp2023api2}. NIST lo
colloca nel gateway o nella mesh anziché nel singolo
servizio~\cite{nist800204}: scelta corretta purché ogni host esposto stia
dietro quel gateway, mentre il limite scritto nel codice viaggia con
l'applicazione in ogni ambiente.

\paragraph{Rilevabilità --- stato dell'arte in letteratura}

Le due fonti primarie sembrano contraddirsi. OWASP classifica la rilevabilità
come facile, poiché richieste con parametri di dimensione crescente, seguite
dall'analisi di stato, tempo e lunghezza della risposta, bastano di norma a
identificare il problema~\cite{owasp2023api4}. CWE-770 giudica invece l'analisi
statica automatica di utilità limitata per le allocazioni senza
limite~\cite{cwe770}. La contraddizione è apparente: OWASP descrive il tester
in black box, CWE lo strumento che analizza il codice. Il test dinamico osserva
il sistema nel suo insieme e vede il limite ovunque sia applicato, ma osserva
solo l'host interrogato, e un host beta privo di gateway resta invisibile
finché non viene inventariato. Sul versante statico la letteratura copre
sottoclassi delimitate, il ReDoS~\cite{staicu2018} e il costo delle query
GraphQL~\cite{cha2020graphql}, mentre manca, per quanto è stato possibile
verificare, un corpo di lavori sottoposti a revisione paritaria sul
rilevamento statico dell'assenza di rate limiting, timeout o limiti di
paginazione nel codice di API REST. Le pratiche esistenti provengono dagli
strumenti e agiscono sul codice, sui manifest infrastrutturali e sul contratto
OpenAPI; solo l'osservazione di CWE-770 sui file di configurazione privi di un
valore massimo legittima le ultime due direzioni a partire da una fonte
primaria~\cite{cwe770}.

Ne discende una conclusione sulla natura del segnale statico che distingue API4
dalle categorie precedenti. Per una BOLA o un difetto di autenticazione
l'assenza del controllo nel gestore coincide con la vulnerabilità, perché il
controllo non può stare altrove. Qui l'analisi del codice non può affermare la
presenza del difetto, dato che il limite potrebbe risiedere nel gateway o nella
piattaforma di orchestrazione; può però affermarne l'assenza a un dato livello,
e tale assenza vale come finding solo se confermata altrove o da una sonda
dinamica. Lo scenario di API9 mostra il caso opposto, in cui la presenza del
limite a un livello non copre tutti gli host~\cite{owasp2023api9}. La
correlazione fra segnali eterogenei è quindi, per questa categoria, condizione
necessaria per formulare un giudizio.

\subsection{API6:2023 -- Unrestricted Access to Sensitive Business Flows}
\label{subsubsec:api6}

\paragraph{Definizione e collocazione}
Le categorie esaminate finora descrivono controlli assenti. Qui ogni singola
chiamata è legittima: l'utente ha diritto di compierla, i controlli di
autorizzazione funzionano, il codice si comporta come progettato. Il danno
nasce dal numero. OWASP definisce vulnerabile l'endpoint che espone un flusso
di business sensibile senza limitarne adeguatamente l'accesso, portando come
esempi l'acquisto di un prodotto, la creazione di un commento e la prenotazione
di una fascia oraria~\cite{owasp_api6_2023}.

Segue l'affermazione che definisce davvero la categoria: il rischio di un
accesso eccessivo cambia da settore a settore, e la creazione di post tramite
script può costituire spam per un social network ed essere incoraggiata da un
altro~\cite{owasp_api6_2023}. Nessuna proprietà del codice, del contratto o
della configurazione rende dunque un endpoint vulnerabile ad API6: lo stesso
identico endpoint è sicuro in un contesto e vulnerabile in un altro, e a
decidere è il modello di business che gli sta dietro.

\paragraph{Cause e meccanismi tecnici tipici}
La causa che OWASP indica non è tecnica: manca una visione d'insieme dell'API
che ne supporti i requisiti di business~\cite{owasp_api6_2023}. Il difetto
nasce nel raccordo fra chi progetta l'endpoint e chi conosce il valore
economico di ciò che l'endpoint muove, e dove quel raccordo manca nessuno è
nella posizione di accorgersene.

Lo sfruttamento richiede di comprendere il modello di business, individuare i
flussi sensibili e automatizzarne l'accesso~\cite{owasp_api6_2023}: soltanto
l'ultimo passo produce traffico osservabile. Il primo scenario lo illustra con
una console da gioco annunciata in quantità limitata, acquistata da un
programma eseguito il giorno del lancio e distribuito su indirizzi IP e località
diverse~\cite{owasp_api6_2023}. La distribuzione geografica non è un dettaglio
di colore: è la contromisura dell'attaccante contro il limite per indirizzo, e
spiega perché il rate limiting nella sua forma più semplice non basti.

\paragraph{Impatto essenziale}
Su questo punto il documento è internamente incoerente. La riga dei rating
classifica l'impatto tecnico come moderato, mentre la descrizione degli impatti
afferma che un impatto tecnico in generale non è atteso e porta come esempi di
danno l'impossibilità per gli utenti legittimi di acquistare un prodotto e
l'inflazione nell'economia interna di un videogioco~\cite{owasp_api6_2023}. È
la seconda formulazione a corrispondere agli scenari descritti, che non
producono né violazione di riservatezza né alterazione di dati né
indisponibilità del servizio. Ciò che cambia sta fuori dal sistema e si misura
in scorte sottratte, ricavi persi o costi di acquisizione pagati per utenti
inesistenti. Ne segue un effetto collaterale rilevante: un incidente che non
produce indisponibilità né perdita di dati non attiva le procedure di risposta
e viene percepito come problema commerciale anziché di sicurezza.

\paragraph{Mitigazione}
OWASP propone una struttura a due livelli in cui l'ordine conta: prima si
identificano, sul piano del business, i flussi che un uso eccessivo può
danneggiare, e solo dopo si sceglie il meccanismo tecnico adeguato al rischio
individuato~\cite{owasp_api6_2023}. Nessuna categoria precedente prescrive una
fase preliminare simile, perché in nessuna altra la configurazione corretta
dipende dall'azienda che la adotta. I meccanismi elencati, dal fingerprinting
del dispositivo al captcha, dall'individuazione di comportamenti non umani come
il passaggio dal carrello all'acquisto in meno di un secondo fino al blocco dei
nodi di uscita Tor, sono presentati come modi per rallentare le minacce
automatizzate, non per impedirle~\cite{owasp_api6_2023}. La logica che li
accomuna è economica: nessuno produce una decisione binaria affidabile, tutti
spostano il costo dell'attacco fino a superare il profitto atteso.

\paragraph{Rilevabilità --- stato dell'arte in letteratura}
Il rilevamento automatico si scontra qui con l'assenza di un oracolo: nessun
artefatto analizzabile, che sia codice, contratto o configurazione, risponde
alla domanda se un dato flusso in una data azienda vada protetto, e la mancanza
di una mappatura CWE registra il fatto che non esiste una debolezza del
software da nominare. Le tecniche note per i difetti di logica applicativa
cercano deviazioni da un comportamento
atteso~\cite{felmetsger2010logic,pellegrino2014logic}, mentre qui il flusso
viene percorso correttamente e l'attacco sta nel ripeterlo. Restano le
contromisure comportamentali a runtime indicate da OWASP~\cite{owasp_api6_2023},
che riconoscono l'automazione durante l'attacco anziché la vulnerabilità prima
che avvenga.

\newpage
\subsection{API7:2023 -- Server Side Request Forgery}
\label{subsubsec:api7}

\paragraph{Definizione e collocazione}
Nelle categorie viste finora l'attaccante agisce sui dati che l'API restituisce o sulle operazioni che esegue. In API7:2023 agisce su un bersaglio diverso: le richieste che l'API emette verso l'esterno. Il difetto si presenta quando un'API recupera una risorsa remota senza validare l'URL fornito dall'utente, e consente all'attaccante di costringere l'applicazione a inviare una richiesta costruita ad arte verso una destinazione inattesa, anche quando essa è protetta da un firewall o da una VPN~\cite{owasp_api7_2023}. Il server diventa un intermediario involontario: agisce per conto di chi lo interroga, ma con i propri privilegi di rete.

MITRE colloca il difetto esattamente in questi termini. CWE-918 \emph{Server-Side Request Forgery} descrive il server che riceve un URL da un componente a monte e ne recupera il contenuto senza accertarsi a sufficienza che la richiesta sia diretta alla destinazione attesa~\cite{cwe918}, ed è figlia diretta di CWE-441 \emph{Unintended Proxy or Intermediary}, che la letteratura di sicurezza conosce come problema del vicario confuso~\cite{cwe918}. La denominazione alternativa registrata dalla stessa voce, XSPA ovvero \emph{Cross Site Port Attack}, conserva traccia del primo uso documentato della tecnica, la scansione di porte interne attraverso un'applicazione web.

Due proprietà della voce CWE meritano attenzione perché la distinguono da quelle incontrate finora. La prima è l'ordinalità: CWE-918 è classificata come \emph{resultant}, ossia una debolezza tipicamente legata alla presenza di altre debolezze~\cite{cwe918}, e ciò riflette il fatto che l'SSRF raramente è il fine dell'attacco quanto piuttosto il mezzo per raggiungerne un altro. La seconda è la persistenza nelle classifiche: la voce compare nelle CWE Top~25 del 2021, 2022, 2023, 2024 e 2025, cinque edizioni consecutive~\cite{cwe918}, e la sua rilevanza è cresciuta al punto che MITRE registra ora fra le tecnologie interessate quelle di intelligenza artificiale con prevalenza frequente, allegando esempi osservati che riguardano toolkit e framework per modelli linguistici~\cite{cwe918}.

Anche la collocazione nelle tassonomie OWASP generaliste è istruttiva. Nella Top~10 del 2021 l'SSRF aveva una categoria propria, A10:2021, mentre nell'edizione 2025 la voce CWE-918 risulta ricondotta ad A01:2025 \emph{Broken Access Control}~\cite{cwe918}. Il riassorbimento non è arbitrario: la conseguenza principale registrata da MITRE sul piano del controllo di accesso è che fornendo URL verso host o porte inattesi l'attaccante fa apparire la richiesta come proveniente dal server, aggirando così controlli quali i firewall che gli impedirebbero di raggiungere quegli indirizzi direttamente~\cite{cwe918}. L'SSRF, in altre parole, non viola un controllo di accesso applicativo ma ne annulla uno di rete.

Ciò che rende la categoria attuale è il contesto in cui le API operano oggi, e OWASP lo articola su due assi. Più comune, perché webhook, recupero di file da URL, single sign-on personalizzati e anteprime dei collegamenti spingono lo sviluppatore ad accedere a risorse esterne indicate dall'utente. Più pericoloso, perché tecnologie come i fornitori cloud, Kubernetes e Docker espongono canali di gestione e controllo su HTTP a percorsi prevedibili e ben noti, che costituiscono un bersaglio facile; a ciò si aggiunge la difficoltà di limitare il traffico in uscita, data la natura interconnessa delle applicazioni moderne~\cite{owasp_api7_2023}. Le metriche assegnate sono sfruttabilità facile, prevalenza comune, rilevabilità facile e impatto tecnico moderato~\cite{owasp_api7_2023}.

Rispetto alle categorie contigue API7 intrattiene tre rapporti che conviene fissare.

Con API10:2023 il legame è di simmetria sul traffico in uscita. Entrambe riguardano richieste che il server emette anziché ricevere, ma differiscono per chi sceglie la destinazione: in API7 la sceglie l'attaccante attraverso un input non validato, in API10 la destinazione è un fornitore legittimo di cui però l'applicazione si fida troppo, non validandone la risposta né seguendone ciecamente le redirezioni~\cite{owasp_api10_2023}. Le due categorie condividono peraltro una contromisura, il divieto di seguire redirezioni automatiche, che OWASP prescrive in entrambe~\cite{owasp_api7_2023,owasp_api10_2023}.

Con API4:2023 il contatto è sul timeout. Un server indotto a contattare un host controllato dall'attaccante può essere trattenuto indefinitamente da una risposta che non arriva, e OWASP indica fra gli impatti possibili dell'SSRF proprio il denial of service e l'uso del server come proxy per mascherare attività malevole~\cite{owasp_api7_2023}. La difesa è la stessa già discussa per il consumo di risorse, ma il difetto primario resta la libertà nella scelta della destinazione.

Con API3:2023 il rapporto riguarda il punto d'ingresso. Il campo che contiene l'URL è una proprietà dell'oggetto in arrivo, e il primo controllo che avrebbe dovuto agire è quello sulla forma e sul valore ammessi per quella proprietà; se il contratto dichiarasse uno schema ristretto e la validazione lo applicasse, l'SSRF non si presenterebbe. La differenza sta nell'esito: in API3 una proprietà scritta indebitamente altera un dato, in API7 altera il comportamento di rete del server. La Tabella~\ref{tab:api7_relazioni} riassume i tre rapporti.

\begin{table}[H]
\centering
\small
\begin{tabular}{p{2.0cm} p{5.6cm} p{5.6cm}}
\hline
\textbf{Categoria} & \textbf{Punto di contatto con API7} & \textbf{Criterio di distinzione} \\
\hline
API3 & L'URL non validato è una proprietà dell'oggetto in ingresso & Esito: alterazione di un dato (API3) contro alterazione del comportamento di rete (API7) \\
API4 & Timeout sulle chiamate in uscita; l'SSRF può causare denial of service & Difetto primario: assenza di limite (API4) contro destinazione scelta dall'attaccante (API7) \\
API10 & Entrambe riguardano il traffico in uscita; comune il divieto di seguire redirezioni & Chi sceglie la destinazione: l'attaccante (API7) contro un fornitore legittimo ma non validato (API10) \\
\hline
\end{tabular}
\caption{Relazioni fra API7:2023 e le categorie contigue della OWASP API Security Top~10 2023.}
\label{tab:api7_relazioni}
\end{table}

\paragraph{Cause e meccanismi tecnici tipici}
La causa immediata è la mancanza o l'inadeguatezza della validazione di un URI fornito dal client, e OWASP osserva che sono proprio i modelli di sviluppo moderni a incoraggiare l'accesso a URI di provenienza esterna~\cite{owasp_api7_2023}. Il difetto nasce quindi da una funzionalità voluta, non da un errore di codifica: nessuno introduce per sbaglio un endpoint che scarica un'immagine da un indirizzo scelto dall'utente.

Il primo scenario del documento mostra la forma elementare. Un social network consente di caricare l'immagine del profilo indicandone l'URL anziché il file, il che innesca una chiamata a \texttt{POST /api/profile/upload\_picture} con il campo \texttt{picture\_url}. L'attaccante vi inserisce \texttt{localhost:8080} e, misurando il tempo di risposta, deduce se la porta sia aperta, avviando così una scansione della rete interna attraverso l'endpoint~\cite{owasp_api7_2023}. Il dettaglio importante è che l'informazione non arriva dal corpo della risposta ma dalla sua latenza: è la forma cieca dell'attacco, in cui l'attaccante non riceve il contenuto della richiesta interna e deve inferirne l'esito da canali laterali.

Il secondo scenario mostra la forma con ritorno, ed è quella economicamente rilevante. Un prodotto di sicurezza consente di configurare webhook verso un SIEM esterno; durante la creazione del canale il back-end invia una richiesta di prova all'URL indicato e ne mostra all'utente la risposta. L'attaccante sostituisce l'URL del SIEM con quello del servizio di metadati dell'ambiente cloud, \texttt{http://169.254.169.254/latest/meta-data/iam/security-credentials/}, e poiché l'applicazione visualizza la risposta della richiesta di prova, ottiene le credenziali dell'ambiente~\cite{owasp_api7_2023}. Qui si compongono i due assi indicati sopra: il webhook è la funzionalità che rende comune il difetto, l'indirizzo di metadati a percorso noto è ciò che lo rende grave.

Una terza famiglia di meccanismi non riguarda l'assenza di validazione ma il suo aggiramento, e OWASP la segnala includendo fra le contromisure l'uso di un parser di URL collaudato per evitare i problemi causati da incoerenze di analisi sintattica~\cite{owasp_api7_2023}. Gli esempi osservati raccolti da MITRE ne documentano il funzionamento: in un caso, la validazione errata del formato decimale di un indirizzo IP consente l'analisi di formati ottali o esadecimali, permettendo l'aggiramento di un meccanismo di protezione contro l'SSRF~\cite{cwe918}. Rientrano nella stessa famiglia le redirezioni seguite automaticamente, per cui un URL inizialmente lecito rimanda a uno interno, e l'uso di schemi diversi da HTTP: MITRE cita esplicitamente \texttt{file://} per accedere a documenti sul sistema e \texttt{gopher://} o \texttt{tftp://}, che possono offrire all'attaccante un controllo maggiore sul contenuto delle richieste~\cite{cwe918}, oltre a un caso in cui una libreria di download segue automaticamente redirezioni verso \texttt{file://} e \texttt{scp://}~\cite{cwe918}. Il punto generale è che la validazione dell'URL è un problema di parsing, e i problemi di parsing si risolvono male con espressioni regolari scritte a mano.

Un'osservazione di Wessels, Koch, Pellegrino e Johns chiude il quadro delle cause. Analizzando 27.078 applicazioni PHP open source, gli autori rilevano che le difese note sono scarsamente adottate e che la stragrande maggioranza dei flussi di dati controllati dall'attaccante da loro individuati risulta vulnerabile, concludendo che gli sviluppatori sono tuttora in larga parte inconsapevoli degli abusi possibili sulle richieste server-side e della necessità di difendersene~\cite{wessels2024ssrf}. La causa più diffusa, in altri termini, non è una difesa mal implementata ma l'assenza di qualunque difesa.

\paragraph{Impatto essenziale}
OWASP elenca come esiti di uno sfruttamento riuscito l'enumerazione dei servizi interni, tipicamente attraverso scansione di porte, la divulgazione di informazioni, l'aggiramento di firewall o altri meccanismi di sicurezza, e in alcuni casi il denial of service o l'uso del server come proxy per occultare attività malevole~\cite{owasp_api7_2023}. CWE-918 organizza gli stessi effetti sulle tre proprietà: lettura di dati applicativi sulla riservatezza, esecuzione di codice o comandi non autorizzati sull'integrità, aggiramento dei meccanismi di protezione sul controllo di accesso~\cite{cwe918}.

L'elemento che distingue l'impatto dell'SSRF è il cambio di perimetro. L'attaccante non ottiene direttamente un dato: ottiene una posizione, quella del server, e da lì raggiunge ciò che dalla propria posizione gli era precluso. Questo spiega perché la voce CWE sia classificata come debolezza risultante e perché il valore dell'attacco dipenda interamente da cosa sia raggiungibile dall'interno. In un'architettura piatta, dove i servizi interni si fidano reciprocamente perché protetti dal perimetro, un singolo endpoint vulnerabile annulla la protezione di tutti; in un'architettura segmentata, lo stesso endpoint consente poco più di una scansione.

Il caso dei servizi di metadati cloud merita una nota, perché è quello che converte l'SSRF da problema di ricognizione a compromissione. Quei servizi rispondono a un indirizzo fisso e ben noto, non richiedono autenticazione perché presuppongono che solo l'istanza possa raggiungerli, ed espongono le credenziali del ruolo associato alla macchina. Un SSRF con ritorno verso quell'indirizzo produce quindi credenziali valide per l'infrastruttura, non un dato applicativo. Che il rischio non sia teorico lo attestano gli esempi osservati registrati da MITRE, fra cui due vulnerabilità SSRF, in un mail server e in una piattaforma cloud, entrambe indicate come sfruttate attivamente secondo il catalogo KEV di CISA~\cite{cwe918}.

\paragraph{Mitigazione}
Le contromisure che OWASP elenca si dispongono su due piani. Sul piano dell'architettura di rete, la raccomandazione è di isolare il meccanismo di recupero delle risorse, osservando che tali funzionalità mirano di norma a recuperare risorse remote e non interne~\cite{owasp_api7_2023}. È la difesa più robusta perché non dipende dalla correttezza della validazione: se il componente che effettua le richieste in uscita non ha rotta verso la rete interna, un URL malevolo non produce nulla.

Sul piano applicativo il documento prescrive l'uso, ove possibile, di liste di elementi ammessi su tre dimensioni distinte: le origini remote da cui gli utenti sono attesi scaricare risorse, gli schemi di URL e le porte, e i tipi di media accettati per una data funzionalità~\cite{owasp_api7_2023}. La tripartizione è significativa perché ciascuna dimensione blocca una famiglia di aggiramenti diversa, e la seconda in particolare chiude gli schemi alternativi discussi sopra. Seguono la disabilitazione delle redirezioni HTTP, l'uso di un parser di URL collaudato e mantenuto, la validazione e sanificazione di tutti i dati forniti dal client, e il divieto di restituire al client le risposte grezze~\cite{owasp_api7_2023}. Quest'ultima misura non impedisce l'attacco ma ne degrada l'esito dalla forma con ritorno a quella cieca, riducendo drasticamente il valore dell'informazione ottenibile.

CWE-918 aggiunge, nell'esempio dimostrativo, la forma più restrittiva della lista di ammessi, ossia il confronto dell'URL con un insieme finito di valori consentiti anziché con un pattern~\cite{cwe918}. È una misura applicabile solo quando le destinazioni legittime sono note in anticipo, il che esclude proprio i casi che rendono la categoria comune: un webhook definito dall'utente non può essere confrontato con una lista.

Va infine registrata un'ammissione che OWASP fa esplicitamente e che nessun'altra categoria contiene: il rischio SSRF non può sempre essere eliminato del tutto, e nella scelta del meccanismo di protezione occorre considerare i rischi e le esigenze di business~\cite{owasp_api7_2023}. La ragione è che la funzionalità e la vulnerabilità coincidono. Un servizio che offre anteprime di collegamenti arbitrari deve, per definizione, effettuare richieste verso indirizzi che l'utente sceglie; ciò che si può fare è confinare l'effetto di quelle richieste, non impedirle.

\paragraph{Rilevabilità: stato dell'arte in letteratura}
Per la prima volta nel corso di questo capitolo le due fonti primarie concordano, e concordano nel senso favorevole all'analisi statica. OWASP classifica la rilevabilità come facile, precisando però che l'analisi delle richieste e delle risposte è necessaria per individuare il problema e che quando la risposta non viene restituita, nella forma cieca, il rilevamento richiede maggiore sforzo e creatività~\cite{owasp_api7_2023}. CWE-918 è ancora più netta: la sezione dei metodi di rilevamento contiene una sola voce, l'analisi statica automatica, descritta come costruzione di un modello di flusso dei dati e di controllo seguita dalla ricerca di pattern potenzialmente vulnerabili che connettono sorgenti, ossia le origini dell'input, a sink, ossia le destinazioni dove il dato interagisce con componenti esterni, e la sua efficacia è dichiarata alta~\cite{cwe918}.

Il contrasto con quanto osservato per le categorie precedenti è netto e merita di essere esplicitato, perché individua la ragione strutturale per cui l'SSRF si presta all'analisi statica mentre l'autorizzazione non vi si presta. Per API5 lo strumento vede l'assenza di un controllo ma non sa se fosse dovuto; per API4 vede l'assenza di un limite ma non sa se esso sia applicato altrove; per API6 non ha nemmeno un artefatto in cui cercare. Per API7 l'oracolo è interno al codice: esiste un flusso di dati da un input controllabile dall'attaccante a una funzione che emette una richiesta di rete, e l'esistenza di quel flusso è una proprietà del programma, verificabile senza conoscere né il contesto di business né le intenzioni del progettista. La domanda che lo strumento deve risolvere non è se qualcosa manchi, ma se due punti del programma siano connessi.

La letteratura recente conferma il quadro ma ne segnala i limiti pratici. Wessels e colleghi applicano un'analisi di flusso dei dati su un grafo di proprietà del codice costruito a partire dal bytecode PHP, individuando gli input controllabili dall'attaccante che raggiungono funzioni note di richiesta server-side, e sottopongono i risultati a validazione manuale~\cite{wessels2024ssrf}: la scelta stessa di prevedere una fase manuale segnala che il segnale automatico, per quanto forte, non è conclusivo. Il lavoro su Artemis affronta il problema specifico dei falsi positivi in PHP, combinando la costruzione di grafi di chiamata espliciti e impliciti, regole che prevengono la propagazione eccessiva della marcatura quando una stringa contaminata non può più produrre un URL sfruttabile, e la potatura tramite analisi della compatibilità delle condizioni di percorso~\cite{artemis2025}. Sul versante dinamico, SSRFuzz combina l'inferenza dinamica della marcatura con il fuzzing, agganciando i sink noti in un ambiente di esecuzione strumentato per restringere lo spazio degli input~\cite{wang2024ssrfuzz}.

Da questi lavori si ricavano i tre problemi che qualunque rilevatore statico di SSRF deve affrontare, e che valgono anche per il modulo descritto nella parte sperimentale. Il primo è la completezza dell'inventario dei sink: le funzioni che emettono richieste di rete sono numerose, appartengono sia alla libreria standard sia a librerie di terze parti, e ometterne una produce un falso negativo silenzioso. Il secondo è la sovra-marcatura: una stringa contaminata concatenata con altre non produce necessariamente un URL controllabile, e trattare come vulnerabile ogni flusso che raggiunge un sink genera rumore. Il terzo è la validazione intermedia: fra la sorgente e il sink può trovarsi un controllo che rende il flusso innocuo, e distinguere una validazione efficace da una aggirabile richiede di ragionare sulla semantica del controllo, non solo sulla sua presenza. È esattamente il punto in cui il rilevamento dell'SSRF smette di essere un problema di raggiungibilità e torna a somigliare agli altri.

Resta la distinzione fra le due forme dell'attacco, che ha una ricaduta diretta sulla strategia di verifica. Per l'SSRF con ritorno la conferma dinamica è immediata: si punta l'endpoint verso un indirizzo controllato e si osserva se la richiesta arriva. Per l'SSRF cieco la conferma richiede un canale laterale, tipicamente una risoluzione DNS o una connessione verso un servizio d'ascolto dedicato, ed è la ragione per cui OWASP parla di sforzo e creatività~\cite{owasp_api7_2023}. In una pipeline che combina segnale statico e sonda dinamica, come quella adottata nel seguito, la conseguenza operativa è che l'assenza di conferma dinamica non deve declassare il finding statico: significa soltanto che il canale di ritorno non è stato osservato.

\subsubsection{API8:2023 -- Security Misconfiguration}
\label{subsubsec:api8}

\paragraph{Definizione e collocazione}
Le categorie viste finora individuano un difetto in un punto preciso: un controllo che manca in un handler, un limite non impostato, un URL non validato. API8:2023 non ha un punto. OWASP dichiara che la configurazione insicura può verificarsi a qualunque livello dello stack dell'API, dal livello di rete a quello applicativo~\cite{owasp_api8_2023}, e l'elenco degli indicatori di vulnerabilità attraversa effettivamente tutti quei livelli: hardening assente in una qualsiasi parte dello stack o permessi mal configurati sui servizi cloud, patch di sicurezza mancanti o sistemi non aggiornati, funzionalità non necessarie abilitate come verbi HTTP o funzioni di logging, discrepanze nel modo in cui i server della catena HTTP elaborano le richieste in arrivo, TLS assente, direttive di sicurezza o di controllo della cache non inviate ai client, politica CORS mancante o impostata male, messaggi di errore che includono tracce dello stack o espongono altre informazioni sensibili~\cite{owasp_api8_2023}.

Sono otto voci che hanno in comune soltanto una cosa: nessuna di esse è un errore di logica applicativa. In ciascun caso il codice fa quello che deve; è l'ambiente in cui gira, o un'impostazione che qualcuno ha lasciato al valore predefinito, a produrre l'esposizione. Questa è la definizione operativa della categoria, e ne spiega tanto l'ampiezza quanto la difficoltà di trattarla come le altre.

La particolarità tassonomica lo conferma nel modo più netto. Alla categoria OWASP associa sette voci CWE, un numero superiore a quello di qualunque altra voce della Top~10~\cite{owasp_api8_2023}. Il problema è che diverse fra esse non sono debolezze. CWE-16 \emph{Configuration} è una categoria, e MITRE ne dichiara la mappatura \emph{prohibited}, ossia proibita, motivandola con il fatto che si tratta di un raggruppamento informale e non di una debolezza in sé, e aggiungendo che l'uso di categorie per la mappatura è sconsigliato dal 2019~\cite{cwe16}. La nota di manutenzione della stessa voce è quanto di più esplicito si possa desiderare: MITRE riconosce che varie fonti vi mappano comunque nonostante la guida contraria, osserva che in questo caso non esistono debolezze CWE chiare da utilizzare, e nota che una formulazione come \emph{Inappropriate Configuration} suonerebbe più conforme allo stile CWE ma continuerebbe a non indicare un comportamento effettivo del prodotto~\cite{cwe16}. L'indicazione operativa che ne segue è di analizzare la debolezza nel comportamento che la configurazione ha determinato, riconducendola per esempio alla discendenza di CWE-284 per il controllo di accesso o a CWE-400 per la gestione delle risorse~\cite{cwe16}.

Ne discende un criterio che vale la pena fissare, perché ha conseguenze su come un difetto di configurazione va classificato. La configurazione non è una debolezza: è il meccanismo attraverso cui una debolezza viene introdotta. Un'impostazione CORS permissiva non è insicura in astratto, lo diventa perché produce un difetto di controllo di accesso; l'assenza di TLS non è una debolezza autonoma, ma la causa di una trasmissione in chiaro di informazioni sensibili, che è CWE-319 ed è una debolezza vera. Delle sette voci elencate da OWASP, quelle utilizzabili per una mappatura sono infatti quelle che nominano un comportamento e non un ambito: CWE-209 sulla generazione di messaggi d'errore contenenti informazioni sensibili, CWE-319 sulla trasmissione in chiaro, CWE-444 sull'interpretazione incoerente delle richieste HTTP, CWE-942 sulla politica cross-domain permissiva verso domini non fidati~\cite{owasp_api8_2023}.

Le metriche assegnate collocano API8 fra le categorie più gravi dell'elenco: sfruttabilità facile, prevalenza diffusa, rilevabilità facile, impatto tecnico severo~\cite{owasp_api8_2023}. È una combinazione che nella Top~10 condivide con poche altre voci, e la motivazione data per il vettore d'attacco chiarisce perché: gli attaccanti cercano difetti non corretti, endpoint comuni, servizi in esecuzione con configurazioni predefinite insicure, file e directory non protetti, e OWASP osserva che si tratta in larga parte di conoscenza pubblica per la quale possono essere disponibili exploit~\cite{owasp_api8_2023}. Non serve comprendere l'applicazione: serve sapere cosa cercare, e quell'informazione è documentata.

Il rapporto con le altre categorie è, per API8, di tipo particolare: non di parentela ma di sovrapposizione parziale ricorrente. Il difetto di configurazione compare come causa alternativa in almeno tre voci già trattate.

Con API4:2023 il contatto è sui limiti dichiarativi. I vincoli su memoria, CPU e processi che OWASP raccomanda di affidare a container e funzioni serverless~\cite{owasp_api4_2023} vivono in file di orchestrazione, e la loro assenza è tanto un difetto di consumo di risorse quanto una configurazione carente nel senso di API8, che chiede espressamente di rivedere i file di orchestrazione e i servizi cloud~\cite{owasp_api8_2023}.

Con API5:2023 il contatto è sul metodo HTTP, già osservato in quella sede. La raccomandazione di API8 di specificare quali verbi siano ammessi per ciascuna API e di disabilitare tutti gli altri~\cite{owasp_api8_2023} previene per via configurativa il caso in cui una rotta risulti protetta in lettura e scoperta in cancellazione. La differenza sta nel livello a cui il rimedio agisce: nell'applicazione per API5, nella catena dei server per API8.

Con API7:2023 il contatto riguarda la politica di uscita. Lo scenario che OWASP porta per API8 lo mostra con chiarezza, poiché l'esecuzione di codice remoto vi diventa possibile per la combinazione di una configurazione predefinita insicura dell'utilità di logging e di una politica di rete in uscita permissiva~\cite{owasp_api8_2023}: la seconda condizione è esattamente quella che rende sfruttabile anche l'SSRF. La Tabella~\ref{tab:api8_relazioni} riassume.

\begin{table}[H]
\centering
\small
\begin{tabular}{p{2.0cm} p{5.6cm} p{5.6cm}}
\hline
\textbf{Categoria} & \textbf{Punto di contatto con API8} & \textbf{Criterio di distinzione} \\
\hline
API4 & Limiti di risorse dichiarati nei file di orchestrazione & Difetto di consumo (API4) contro impostazione mancante nello stack (API8) \\
API5 & Verbi HTTP non necessari abilitati sulla catena dei server & Livello del rimedio: applicazione (API5) contro configurazione dei server (API8) \\
API7 & Politica di rete in uscita permissiva come precondizione dello sfruttamento & Origine della richiesta: input non validato (API7) contro assenza di segmentazione (API8) \\
\hline
\end{tabular}
\caption{Sovrapposizioni fra API8:2023 e le categorie contigue della OWASP API Security Top~10 2023.}
\label{tab:api8_relazioni}
\end{table}

\paragraph{Cause e meccanismi tecnici tipici}
La causa strutturale è che la configurazione di un'API moderna non risiede in un luogo unico né sotto una responsabilità unica. Si distribuisce fra il codice applicativo, i file di orchestrazione, il reverse proxy, il gateway, le impostazioni del fornitore cloud e i valori predefiniti delle librerie di terze parti, e ciascuno di questi livelli ha un proprietario diverso e un ciclo di vita proprio. Un'impostazione corretta al momento del rilascio può cessare di esserlo dopo un aggiornamento che nessuno ha collegato alla sicurezza.

Il primo scenario di OWASP illustra il meccanismo dei valori predefiniti pericolosi, e lo fa con quello che è riconoscibilmente il caso Log4Shell. Un server di back-end tiene un registro degli accessi scritto da una diffusa utilità di logging open source che supporta l'espansione di segnaposto e le ricerche JNDI, entrambe abilitate per impostazione predefinita; per ogni richiesta viene scritta una riga secondo un formato che include il percorso e il codice di stato. L'attaccante invia una richiesta apparentemente innocua verso \texttt{/health} con un'intestazione \texttt{X-Api-Version} contenente un riferimento JNDI a un server LDAP che controlla, e nel momento in cui l'utilità espande quel valore per scriverlo nel registro scarica ed esegue l'oggetto malevolo dal server remoto~\cite{owasp_api8_2023}. Vale la pena osservare cosa manca in questa catena: non c'è alcun difetto nel codice dell'applicazione, non c'è validazione dell'input che avrebbe potuto essere presente, e l'endpoint colpito è quello di verifica dello stato di salute, tipicamente il più innocuo dell'intera API.

Il secondo scenario mostra il meccanismo opposto, quello dell'impostazione assente. Un social network espone i messaggi privati tramite una chiamata a \texttt{/dm/user\_updates.json} con identificativo di conversazione e cursore; poiché la risposta non include l'intestazione \texttt{Cache-Control}, le conversazioni private finiscono nella cache del browser, da cui possono essere recuperate leggendo i file sul filesystem~\cite{owasp_api8_2023}. Qui non c'è nulla di attivo da disabilitare: c'è una direttiva che nessuno ha inviato, e la conseguenza è che dati riservati vengono conservati in chiaro su una macchina che l'applicazione non controlla.

Un terzo meccanismo, che gli scenari non illustrano ma che l'elenco degli indicatori nomina, riguarda le discrepanze nell'elaborazione delle richieste lungo la catena dei server~\cite{owasp_api8_2023}. È il fenomeno che la voce CWE-444 chiama interpretazione incoerente delle richieste HTTP e che la pratica conosce come \emph{request smuggling}: quando un bilanciatore di carico e un server applicativo delimitano diversamente il confine fra due richieste consecutive sulla stessa connessione, un attaccante può far sì che parte del proprio traffico venga attribuita alla richiesta di un altro utente. Il difetto non risiede in nessuno dei due componenti presi singolarmente, ma nella loro combinazione, ed è la ragione per cui OWASP raccomanda di assicurarsi che tutti i server della catena elaborino le richieste in maniera uniforme~\cite{owasp_api8_2023}.

Resta la causa più banale e più diffusa, quella che l'elenco degli indicatori mette al secondo posto: patch mancanti e sistemi non aggiornati~\cite{owasp_api8_2023}. Non è un errore di configurazione in senso stretto ma di manutenzione, e la sua presenza nella stessa categoria segnala che il confine fra le due cose, nella pratica operativa, è poco netto.

\paragraph{Impatto essenziale}
La formulazione di OWASP è breve e va letta nei suoi due termini: le configurazioni insicure non solo espongono dati sensibili degli utenti, ma anche dettagli del sistema che possono condurre alla compromissione completa del server~\cite{owasp_api8_2023}. La seconda parte è quella che giustifica la classificazione dell'impatto tecnico come severo, e i due scenari ne mostrano i due estremi: nel secondo il danno è l'esposizione di conversazioni private, nel primo è l'esecuzione di codice arbitrario sul server.

Ciò che distingue l'impatto di API8 è la sua natura abilitante. Un difetto di configurazione raramente costituisce l'attacco: costituisce la condizione che rende praticabile un altro attacco, o che ne amplifica l'esito. Una traccia di stack in un messaggio d'errore non compromette nulla di per sé, ma rivela versioni, percorsi interni e talvolta frammenti di query, riducendo il costo della ricognizione; una politica di uscita permissiva non è sfruttabile direttamente, ma trasforma un SSRF cieco in un SSRF con ritorno. Questo spiega perché la categoria compaia come precondizione in tante altre voci dell'elenco e perché il suo impatto sia difficile da valutare in isolamento.

Un secondo tratto riguarda l'ampiezza. Poiché la configurazione è per sua natura trasversale, un singolo difetto non riguarda un endpoint ma tutti gli endpoint serviti da quel componente. L'intestazione di cache mancante nel secondo scenario non affligge una conversazione: affligge ogni risposta che attraversa quella configurazione.

\paragraph{Mitigazione}
La struttura che OWASP propone distingue due piani. Sul piano del processo, il ciclo di vita dell'API dovrebbe comprendere un procedimento di hardening ripetibile che consenta un rilascio rapido e agevole di un ambiente adeguatamente irrigidito, un'attività di revisione e aggiornamento delle configurazioni sull'intero stack che includa i file di orchestrazione, i componenti dell'API e i servizi cloud, e un processo automatizzato che valuti in modo continuo l'efficacia di configurazioni e impostazioni in tutti gli ambienti~\cite{owasp_api8_2023}. I tre requisiti hanno un'implicazione comune che merita di essere resa esplicita: la configurazione va trattata come artefatto versionato e verificabile, non come stato accumulato di una macchina.

Sul piano dei controlli specifici, il documento elenca sei misure~\cite{owasp_api8_2023}. La comunicazione fra client e server e verso ogni componente a monte o a valle deve avvenire su canale cifrato, indipendentemente dal fatto che l'API sia interna o esposta al pubblico, precisazione che chiude l'assunto per cui il traffico interno sarebbe fidato. I verbi HTTP ammessi vanno specificati per ciascuna API e tutti gli altri disabilitati. Le API destinate a client basati su browser devono implementare una politica CORS adeguata e includere le intestazioni di sicurezza applicabili. I tipi di contenuto e i formati accettati in ingresso vanno ristretti a quelli che soddisfano i requisiti funzionali. Tutti i server della catena devono elaborare le richieste in modo uniforme per evitare problemi di desincronizzazione. Infine, dove applicabile, vanno definiti e applicati gli schemi di tutti i payload di risposta, comprese le risposte di errore, per impedire che tracce di eccezione e altre informazioni di valore raggiungano l'attaccante.

Quest'ultima misura merita una nota perché collega la mitigazione al contratto. Uno schema di risposta dichiarato per le condizioni di errore non è documentazione: è un vincolo che, se applicato a runtime, impedisce strutturalmente la fuga di informazioni diagnostiche. È il caso in cui la specifica dell'API funge da controllo di sicurezza e non da suo commento.

\paragraph{Rilevabilità: stato dell'arte in letteratura}
OWASP classifica la rilevabilità come facile e ne dà una motivazione che, per la prima volta in questo capitolo, rinvia direttamente agli strumenti: sono disponibili strumenti automatici per individuare e sfruttare configurazioni errate come servizi non necessari o opzioni obsolete~\cite{owasp_api8_2023}. La formulazione è simmetrica e va notata: gli stessi strumenti servono a chi difende e a chi attacca, perché ciò che cercano è pubblicamente noto.

La ragione per cui il rilevamento funziona bene si comprende per contrasto con le categorie precedenti. Per l'autorizzazione funzionale lo strumento non sapeva se un controllo fosse dovuto; per il consumo di risorse non sapeva se un limite fosse applicato altrove; per i flussi di business non aveva alcun artefatto da interrogare. Per la configurazione l'oracolo è esterno e pubblico: esistono valori raccomandati, elenchi di impostazioni predefinite pericolose, versioni note come vulnerabili, e verificare la conformità di una configurazione rispetto a un riferimento è un problema di confronto, non di inferenza. È il motivo per cui questa è l'unica categoria della Top~10 per cui esiste un mercato maturo di strumenti che operano su artefatti dichiarativi anziché su codice.

Il quadro presenta però tre limiti che vanno riconosciuti.

Il primo riguarda la copertura degli artefatti. La configurazione risiede in file di orchestrazione, manifest di deployment, definizioni infrastrutturali, impostazioni del gateway, valori predefiniti delle librerie e stato effettivo delle risorse cloud; nessuno strumento li legge tutti, e ciò che non è dichiarato in un file non è ispezionabile prima dell'esecuzione. OWASP stessa, chiedendo un processo automatizzato che valuti l'efficacia delle impostazioni in tutti gli ambienti~\cite{owasp_api8_2023}, riconosce implicitamente che la verifica sugli artefatti non basta e va accompagnata dall'osservazione dell'ambiente reale.

Il secondo riguarda la deriva. Un'analisi statica di un file infrastrutturale attesta ciò che quel file dichiara, non ciò che è effettivamente attivo sul sistema; ogni modifica applicata fuori dalla pipeline, e ogni impostazione modificata manualmente in produzione, apre una divergenza che l'analisi non vede. Il difetto non è dello strumento ma della premessa: la configurazione descritta e quella effettiva sono due oggetti distinti.

Il terzo riguarda il contesto. Una politica CORS che ammette qualunque origine è un difetto per un'API destinata a client autenticati e una scelta legittima per un'API pubblica di sola lettura; un'intestazione di cache assente è irrilevante su una risposta pubblica e grave su una privata, come mostra il secondo scenario OWASP. La verifica sintattica di un'impostazione, in altri termini, non stabilisce se essa sia sbagliata: stabilisce se si discosta da un riferimento, e la pertinenza di quel riferimento dipende dalla funzione dell'endpoint. È una forma attenuata della difficoltà già incontrata per i flussi di business, con la differenza che qui il contesto necessario è tecnico e ricavabile in parte dal contratto, non economico e del tutto esterno agli artefatti.

Resta un'osservazione sulla natura del segnale. Per le categorie di autorizzazione l'analisi cercava l'assenza di qualcosa che avrebbe dovuto esserci; per API8 cerca la presenza di qualcosa che non avrebbe dovuto esserci, come un verbo abilitato o una direttiva permissiva, oppure l'assenza di qualcosa che ha un valore corretto noto, come un'intestazione di sicurezza. Sono due problemi asimmetrici: il secondo si presta a una verifica per confronto con un riferimento e produce falsi positivi contestuali ma pochi falsi negativi, mentre il primo richiede di ricostruire l'intenzione del progettista. Questa asimmetria spiega perché la stessa infrastruttura di analisi statica dia risultati di qualità molto diversa a seconda della categoria a cui viene applicata.

hiarate sono candidate a essere API ombra, cioè funzionalità raggiungibili che nessun documento menziona e che quindi sfuggono a ogni processo di revisione basato sul contratto. Le rotte dichiarate ma non implementate sono, all'opposto, un indizio di documentazione obsoleta, e possono segnalare una dismissione avvenuta senza aggiornare la specifica oppure una specifica scritta in anticipo e mai realizzata. Nessuna delle due differenze è di per sé una vulnerabilità, ma entrambe misurano la divergenza fra il modello e la realtà, che è esattamente il difetto che la categoria descrive.

Questo confronto ha però un limite intrinseco che va dichiarato: opera fra due rappresentazioni della stessa applicazione, e non dice nulla sugli host su cui quell'applicazione è distribuita. Il primo scenario di OWASP è precisamente il caso in cui codice e contratto coincidono perfettamente e il difetto risiede altrove, nella topologia di rete. Rilevare quel caso richiede di enumerare gli host effettivamente raggiungibili e di confrontarli con l'inventario dichiarato, che è un'attività di ricognizione infrastrutturale e non di analisi del software.

La letteratura accademica su questo problema specifico è scarsa, e vale la pena dirlo esplicitamente anziché colmare il vuoto con riferimenti impropri. La ragione plausibile è che il problema non si presta alla forma sperimentale consueta della ricerca sulla sicurezza del software: non esiste un insieme di verità di riferimento contro cui misurare un rilevatore di API ombra, perché costruire quell'insieme richiederebbe di conoscere l'inventario completo di organizzazioni reali, che è precisamente l'informazione che per ipotesi non esiste. Il campo è quindi presidiato da strumenti commerciali di scoperta delle API, che operano tipicamente osservando il traffico di produzione, e le cui prestazioni sono documentate da materiale dei fornitori e non da valutazioni indipendenti sottoposte a revisione paritaria.

Ne discende una conclusione sulla natura del segnale che chiude la trattazione di questa categoria. Per API9 non esiste un rilevamento nel senso in cui esiste per le categorie precedenti, cioè l'individuazione di un difetto in un artefatto. Esiste la misura di una divergenza fra rappresentazioni indipendenti dello stesso sistema — il contratto, il codice, l'inventario degli host, il traffico osservato — e ogni divergenza è un indizio la cui gravità dipende da cosa si trovi dall'altra parte. Il valore di questa misura non sta nel produrre un verdetto ma nel restituire visibilità: ciò che il difetto sottrae per definizione.


\subsubsection{API10:2023 -- Unsafe Consumption of APIs}
\label{subsubsec:api10}

\paragraph{Definizione e collocazione}
Le nove categorie esaminate finora hanno un tratto comune che diventa visibile solo quando si incontra la decima: descrivono tutte l'API come bersaglio di un attacco diretto. API10:2023 rovescia la prospettiva. Qui l'attaccante non colpisce l'API bersaglio ma un servizio da cui essa dipende, e il danno arriva attraverso un canale che l'applicazione considera fidato. Le note di rilascio dell'edizione 2023 dichiarano di aver introdotto la categoria per affrontare un fenomeno osservato di recente, ossia il fatto che gli attaccanti hanno cominciato a cercare i servizi integrati di un bersaglio per compromettere quelli, anziché colpire direttamente le API del bersaglio stesso~\cite{owasp_api_2023_release_notes}.

La causa che OWASP individua è di ordine psicologico prima che tecnico, e il documento la enuncia senza mediazioni: gli sviluppatori tendono a fidarsi dei dati ricevuti da API di terze parti più di quanto si fidino dell'input dell'utente, e questo vale in particolare per le API offerte da aziende note; di conseguenza tendono ad adottare standard di sicurezza più deboli, per esempio in materia di validazione e sanificazione dell'input~\cite{owasp_api10_2023}. La stessa osservazione ritorna nella descrizione della debolezza, dove si legge che gli sviluppatori tendono a fidarsi senza verificare degli endpoint che interagiscono con API esterne o di terze parti, affidandosi a requisiti più deboli riguardo alla sicurezza del trasporto, all'autenticazione e autorizzazione, e alla validazione e sanificazione dell'input~\cite{owasp_api10_2023}.

Il documento elenca cinque condizioni che rendono l'API vulnerabile: interagire con altre API su un canale non cifrato; non validare e sanificare adeguatamente i dati raccolti da altre API prima di elaborarli o di passarli a componenti a valle; seguire ciecamente le redirezioni; non limitare la quantità di risorse disponibili per elaborare le risposte dei servizi di terze parti; non implementare timeout per le interazioni con essi~\cite{owasp_api10_2023}.

Su questo elenco vale la pena registrare un'obiezione sollevata pubblicamente sul repository del progetto durante la fase di revisione, secondo cui almeno tre di quelle voci apparterrebbero concettualmente ad altre categorie: il limite alle risorse per elaborare le risposte ricade nel consumo incontrollato di risorse, mentre l'interazione su canale non cifrato ricade nella configurazione insicura~\cite{owasp_api_security_issue80}. L'osservazione è fondata e il gruppo di lavoro non l'ha recepita, sicché l'elenco pubblicato mescola condizioni che appartengono al nucleo della categoria, come la mancata validazione dei dati ricevuti e le redirezioni seguite ciecamente, con altre che sono la manifestazione di difetti già trattati altrove nella direzione del traffico in uscita.

I riferimenti esterni indicati sono tre — CWE-20 sulla validazione impropria dell'input, CWE-200 sull'esposizione di informazioni sensibili a un attore non autorizzato, CWE-319 sulla trasmissione in chiaro~\cite{owasp_api10_2023} — e la loro genericità riflette una proprietà reale della categoria anziché una scelta approssimativa. OWASP scrive infatti che l'impatto varia a seconda di ciò che l'API bersaglio fa con i dati che ha prelevato~\cite{owasp_api10_2023}: la debolezza specifica dipende dalla destinazione del dato non validato, che può essere una query SQL, un motore di template o un componente a valle qualsiasi. API10 individua quindi un punto d'ingresso, non un difetto terminale.

Le metriche assegnate presentano una tensione interna che merita di essere segnalata. La sfruttabilità è dichiarata facile, ma la descrizione del vettore afferma che lo sfruttamento richiede all'attaccante di identificare e potenzialmente compromettere altre API o servizi con cui l'API bersaglio è integrata, aggiungendo che di norma tale informazione non è pubblicamente disponibile o il servizio integrato non è facilmente attaccabile~\cite{owasp_api10_2023}. La prevalenza è comune, la rilevabilità media e l'impatto tecnico severo~\cite{owasp_api10_2023}. Se ne ricava che il giudizio di facilità riguarda la fase finale dell'attacco, quella in cui il dato malevolo viene consegnato e accettato senza controlli, mentre la fase preparatoria resta la parte costosa.

Il rapporto con le categorie contigue è definito dalla direzione del traffico e dalla natura della fiducia mal riposta.

Con API7:2023 la simmetria è quella già osservata in quella sede: entrambe riguardano richieste che il server emette, ma in API7 la destinazione è scelta dall'attaccante attraverso un input non validato, mentre in API10 la destinazione è legittima e il problema riguarda ciò che torna indietro. Le due categorie condividono la contromisura del divieto di seguire redirezioni automatiche~\cite{owasp_api7_2023,owasp_api10_2023}.

Con API4:2023 e con API8:2023 i rapporti sono quelli evidenziati dall'obiezione riportata sopra: i timeout e i limiti sulle risposte appartengono al consumo di risorse, la cifratura del canale alla configurazione. In entrambi i casi la specificità di API10 non sta nel controllo mancante ma nella sua destinazione, cioè nel fatto che il componente non protetto sia rivolto verso un fornitore anziché verso un client.

Con API9:2023 il legame è complementare e riguarda il medesimo oggetto visto da due lati. Il punto cieco dei flussi di dati descritto in quella categoria concerne ciò che l'organizzazione condivide con terze parti senza saperlo; API10 concerne ciò che riceve indietro senza verificarlo. Non a caso la prima raccomandazione di API9 sul secondo inventario chiede di censire i servizi integrati documentandone ruolo, dati scambiati e sensibilità~\cite{owasp_api9_2023}: quell'inventario è la precondizione di qualunque valutazione della postura di sicurezza dei fornitori, che è a sua volta la prima contromisura prescritta da API10. La Tabella~\ref{tab:api10_relazioni} riassume.

\begin{table}[H]
\centering
\small
\begin{tabular}{p{2.0cm} p{5.6cm} p{5.6cm}}
\hline
\textbf{Categoria} & \textbf{Punto di contatto con API10} & \textbf{Criterio di distinzione} \\
\hline
API4 & Timeout e limiti sulle risposte dei servizi integrati & Bene protetto: le risorse del server (API4) contro l'integrità del dato ricevuto (API10) \\
API7 & Entrambe riguardano il traffico in uscita e vietano le redirezioni cieche & Destinazione scelta dall'attaccante (API7) contro fornitore legittimo ma non verificato (API10) \\
API8 & Interazione con servizi terzi su canale non cifrato & Impostazione dello stack (API8) contro fiducia mal riposta nella dipendenza (API10) \\
API9 & Inventario dei servizi integrati e dei dati scambiati & Cosa si condivide senza saperlo (API9) contro cosa si accetta senza verificarlo (API10) \\
\hline
\end{tabular}
\caption{Relazioni fra API10:2023 e le categorie contigue della OWASP API Security Top~10 2023.}
\label{tab:api10_relazioni}
\end{table}

\paragraph{Cause e meccanismi tecnici tipici}
La causa strutturale è un confine di fiducia collocato nel punto sbagliato. L'architettura difensiva di un'applicazione presuppone di norma che esista un perimetro, con l'input dell'utente all'esterno e i componenti interni all'interno, e concentra i controlli sull'attraversamento di quel confine. Un servizio di terze parti si trova, in questo schema, in una posizione ambigua: è esterno all'applicazione ma è stato scelto deliberatamente, è raggiungibile su un canale che l'applicazione stessa ha configurato, e appartiene spesso a un'azienda la cui reputazione induce a considerarlo affidabile. La conseguenza è che la risposta del fornitore entra nel sistema attraversando un confine che nessuno ha presidiato, perché nessuno lo ha riconosciuto come confine.

Il primo scenario mostra il meccanismo nella sua forma più pura, quella dell'iniezione differita. Un'API si affida a un servizio di terze parti per arricchire gli indirizzi commerciali forniti dagli utenti: l'indirizzo viene inviato al servizio esterno e i dati restituiti sono poi salvati in una base di dati SQL locale. Gli attaccanti usano il servizio di terze parti per memorizzare un payload di SQL injection associato a un'attività commerciale creata da loro; poi si rivolgono all'API vulnerabile fornendo un input specifico che le fa prelevare la loro attività malevola dal servizio esterno, e il payload finisce per essere eseguito dalla base di dati, esfiltrando dati verso un server controllato dall'attaccante~\cite{owasp_api10_2023}.

Conviene isolare cosa rende questo scenario diverso da una normale iniezione. L'attaccante non invia mai il payload all'API bersaglio: lo deposita altrove, in un sistema che l'API bersaglio interroga per ragioni proprie. Se la difesa dell'applicazione consiste nel sanificare l'input degli utenti, quella difesa non viene mai attivata, perché l'input dell'utente in questo caso è un innocuo nome di attività commerciale. Il payload arriva dal canale che l'applicazione considera fidato, e arriva in un momento diverso da quello della richiesta che lo ha provocato.

Il secondo scenario mostra il meccanismo della redirezione seguita ciecamente. Un'API si integra con un fornitore terzo per conservare in modo sicuro informazioni mediche sensibili, inviando i dati su connessione cifrata tramite una richiesta a \texttt{POST /user/store\_phr\_record} il cui corpo contiene una sequenza genomica. Gli attaccanti trovano il modo di compromettere l'API del fornitore, che comincia a rispondere a quelle richieste con un \texttt{308 Permanent Redirect} verso \texttt{attacker.com}; poiché l'API segue ciecamente le redirezioni del fornitore, ripete la richiesta identica, dati sensibili dell'utente compresi, verso il server dell'attaccante~\cite{owasp_api10_2023}. Il dettaglio decisivo è che qui tutto è stato fatto correttamente sul piano della cifratura e della scelta del fornitore: la fuga avviene perché il client HTTP applica un comportamento predefinito, il seguire le redirezioni, che nessuno ha valutato come decisione di sicurezza.

Il terzo scenario è il più breve e il più istruttivo, perché mostra che la terza parte non deve essere compromessa perché l'attacco funzioni. Un attaccante prepara un repository Git denominato \texttt{'; drop db;--}; quando un'applicazione bersaglio esegue un'integrazione con quel repository, il payload di SQL injection viene usato da un'applicazione che costruisce una query SQL ritenendo che il nome del repository sia un input sicuro~\cite{owasp_api10_2023}. Qui il servizio integrato si comporta in modo del tutto corretto: restituisce fedelmente il nome che l'utente ha scelto per il proprio repository. È l'applicazione consumatrice ad aver confuso l'autenticità del canale con l'affidabilità del contenuto, due proprietà che non hanno alcun rapporto fra loro.

Da questi tre casi si ricava la generalizzazione utile. Un servizio di terze parti può veicolare dati malevoli in tre modi indipendenti: perché è stato compromesso, perché è stato usato come deposito da un attaccante che ne è utente legittimo, o perché trasmette fedelmente contenuti che utenti terzi vi hanno immesso. Solo il primo caso richiede un'intrusione, e i restanti due sono alla portata di chiunque possa registrarsi al servizio.

\paragraph{Impatto essenziale}
OWASP dichiara che l'impatto varia a seconda di ciò che l'API bersaglio fa con i dati prelevati, e che uno sfruttamento riuscito può portare all'esposizione di informazioni sensibili ad attori non autorizzati, a molti tipi di iniezione, o a denial of service~\cite{owasp_api10_2023}. L'impatto tecnico è classificato come severo, e i tre scenari coprono effettivamente tre punti diversi dello spettro: esfiltrazione tramite iniezione nel primo, fuga diretta di dati sanitari nel secondo, distruzione di dati nel terzo.

Ciò che caratterizza l'impatto di questa categoria è che il difetto non ne determina l'entità. Nelle categorie di autorizzazione la gravità è funzione di ciò che il controllo mancante proteggeva; qui è funzione di dove il dato non validato finisce, e quel dato può attraversare diversi componenti prima di arrivare a destinazione. Un'informazione prelevata da un fornitore e memorizzata può essere riletta mesi dopo da un componente diverso, scritto da persone diverse, che la tratta come dato interno perché proviene dalla base di dati dell'organizzazione. La contaminazione, in altri termini, sopravvive alla richiesta che l'ha introdotta.

Un secondo tratto riguarda la scala del danno potenziale. Poiché il vettore passa per una dipendenza condivisa, la compromissione di un singolo fornitore raggiunge simultaneamente tutte le applicazioni che lo integrano, e ciascuna di esse subisce l'attacco senza che nulla nel proprio perimetro sia stato violato. È la ragione per cui OWASP colloca la categoria nel discorso sugli attacchi alla catena di fornitura, e per cui la valutazione del rischio di un'integrazione non può limitarsi alla funzionalità che essa fornisce.

\paragraph{Mitigazione}
Le contromisure che OWASP elenca sono quattro e coprono piani distinti~\cite{owasp_api10_2023}.

La prima è organizzativa e precede lo sviluppo: nel valutare i fornitori di servizi, occorre esaminarne la postura di sicurezza. È una raccomandazione che si applica al momento della scelta e che richiede di trattare l'integrazione come una decisione architetturale con implicazioni di sicurezza, non come una dipendenza tecnica da aggiungere.

La seconda riguarda il trasporto: tutte le interazioni con l'API devono avvenire su un canale di comunicazione sicuro. È la misura che risponde alla prima delle cinque condizioni di vulnerabilità e che, come osservato, appartiene propriamente alla configurazione dello stack.

La terza è il nucleo della categoria: i dati ricevuti dalle API integrate vanno sempre validati e adeguatamente sanificati prima di essere usati. La formulazione è deliberatamente identica a quella che si applicherebbe all'input dell'utente, ed è proprio questo il punto: la contromisura consiste nell'abolire la distinzione fra le due sorgenti.

La quarta riguarda le redirezioni: va mantenuta una lista di posizioni note verso cui le API integrate possono redirigere quelle proprie, e le redirezioni non vanno seguite ciecamente. Vale la pena notare la forma della prescrizione, che non è un divieto assoluto ma una lista di destinazioni ammesse, riconoscendo che alcune redirezioni sono legittime nel funzionamento ordinario dei servizi.

A queste quattro misure il documento non ne aggiunge altre, ma le cinque condizioni di vulnerabilità ne implicano una quinta, relativa ai timeout e ai limiti sulle risposte, che l'elenco delle prevenzioni non riprende. È un'asimmetria interna al documento: due delle condizioni enunciate non trovano una contromisura corrispondente, e vanno cercate nella trattazione del consumo di risorse.

Una misura ulteriore, che discende dalla natura del difetto piuttosto che dal testo, riguarda la forma della risposta attesa. Se il contratto del fornitore è noto, la risposta può essere validata rispetto a uno schema che ne vincoli tipi, lunghezze e formati, esattamente come si fa per le richieste in ingresso. Ciò non intercetta un payload che rispetti lo schema, come il nome di repository del terzo scenario, ma riduce la superficie a ciò che è sintatticamente ammissibile e trasforma la fiducia implicita in un vincolo dichiarato.

\paragraph{Rilevabilità: stato dell'arte in letteratura}
OWASP classifica la rilevabilità come media e non fornisce, a differenza di altre categorie, indicazioni su come procedere; la descrizione della debolezza si limita a osservare che gli attaccanti devono identificare i servizi con cui il bersaglio si integra, cioè le sorgenti dei dati, ed eventualmente comprometterli~\cite{owasp_api10_2023}. L'assenza di una metodologia di verifica è essa stessa un'informazione, e conviene ricostruire il motivo.

Sul piano dell'analisi statica il problema si formula in termini precisi, e la formulazione rivela perché gli strumenti esistenti raramente lo intercettano. L'analisi della propagazione della contaminazione, quella su cui si fonda il rilevamento delle iniezioni e dell'SSRF, richiede di specificare un insieme di sorgenti e un insieme di destinazioni; per convenzione consolidata le sorgenti sono i punti in cui entra l'input dell'utente, ossia parametri della richiesta, corpo, intestazioni e cookie. Per API10 la sorgente è un'altra: è il valore restituito da una chiamata HTTP in uscita. Un motore configurato secondo la convenzione ordinaria non considera contaminato quel valore, e di conseguenza un flusso che va dalla risposta di un fornitore a una query SQL non produce alcuna segnalazione, benché sia strutturalmente identico a un flusso che produrrebbe una segnalazione se partisse da un parametro di richiesta.

Ne discende che il rilevamento di questa categoria non richiede una tecnica nuova ma una diversa configurazione di una tecnica nota: l'estensione dell'insieme delle sorgenti alle risposte delle chiamate in uscita. Il costo di questa estensione è però prevedibile, ed è l'aumento dei falsi positivi, poiché la maggior parte dei dati che un'applicazione preleva da servizi esterni viene usata in modi del tutto innocui. Distinguere i flussi rilevanti richiede di sapere quali destinazioni siano sensibili, il che riporta al problema generale della precisione nell'analisi della contaminazione, e di sapere se fra sorgente e destinazione sia interposta una validazione efficace, che è il problema già incontrato nella trattazione dell'SSRF.

Restano due difetti fra i cinque enunciati che si prestano a un rilevamento più diretto, perché non richiedono di seguire un flusso ma di osservare una configurazione locale della chiamata. Il primo è la redirezione seguita per impostazione predefinita: la maggior parte dei client HTTP la abilita senza che lo sviluppatore debba dichiararlo, sicché il segnale non è la presenza di un'opzione errata ma l'assenza di un'opzione che disattivi il comportamento. Il secondo è il timeout mancante, che ha la stessa struttura. In entrambi i casi si tratta di verificare che una chiamata in uscita dichiari esplicitamente una scelta, e l'assenza della dichiarazione equivale all'accettazione del valore predefinito, che è quello insicuro.

Sul piano dinamico la verifica è ostacolata da un problema di controllo. Provare che un'applicazione accetti dati malevoli da un fornitore richiede di fare in modo che il fornitore li restituisca, il che è possibile solo se il servizio consente all'attaccante di immettervi contenuti, come nel primo e nel terzo scenario, oppure sostituendo il fornitore con un doppio controllato durante il collaudo. La seconda strada è praticabile in un ambiente di prova ma verifica il comportamento dell'applicazione, non la sicurezza dell'integrazione reale; la prima dipende dalle caratteristiche del servizio specifico e non si generalizza.

Va infine registrato che la letteratura sottoposta a revisione paritaria dedicata specificamente a questa categoria è pressoché assente, per quanto è stato possibile verificare. Esiste un corpo consistente di lavori sugli attacchi alla catena di fornitura del software, che però riguardano prevalentemente le dipendenze da pacchetti e la loro compromissione, e un corpo altrettanto consistente sull'analisi della contaminazione, che assume la convenzione sulle sorgenti discussa sopra. Il caso specifico dell'API che si fida della risposta di un'altra API si colloca fra i due e non è stato, finora, oggetto di trattazione autonoma. La sua introduzione nella Top~10 soltanto nell'edizione 2023, motivata dall'osservazione di un fenomeno emergente~\cite{owasp_api_2023_release_notes}, rende plausibile che il ritardo della letteratura rifletta il ritardo con cui il problema è stato riconosciuto come classe a sé.
